% 版式沿用用户样张《粒度与重点样张.tex》：ctexart + 蓝色结论框 + 红色局部重点框。
% XeLaTeX 编译。内容依据标准教材与课程知识整理，非用户资料抽取。
\documentclass[11pt,a4paper]{ctexart}
\usepackage[margin=2cm]{geometry}
\usepackage{amsmath,amssymb,mathtools,bm}
\usepackage[most]{tcolorbox}
\usepackage[hidelinks]{hyperref}
\definecolor{jl}{RGB}{20,70,160}
\definecolor{zd}{RGB}{210,45,55}
\newtcolorbox{结论}[1]{enhanced,colback=jl!3,colframe=jl,
 boxrule=.7pt,arc=1pt,left=7pt,right=7pt,top=5pt,bottom=5pt,
 before skip=9pt,after skip=9pt,fonttitle=\bfseries,title={#1}}
\newcommand{\重点}[1]{\begingroup\setlength{\fboxsep}{3pt}%
 \fcolorbox{zd}{white}{\ensuremath{\displaystyle #1}}\endgroup}
% 文字版重点框：参数是正文（可含 $...$ 和中文）。
\newcommand{\重点文}[1]{\begingroup\setlength{\fboxsep}{2pt}\fcolorbox{zd}{white}{#1}\endgroup}
% 警告与重点区分：警告用文字标“注意”，不另用红框。
\newcommand{\注意}{\textbf{\color{zd}注意}：}
\newcommand{\参考}[1]{\par\noindent\begin{minipage}{\linewidth}\small\color{gray}参考：#1\end{minipage}\par}
\newcommand{\dd}{\mathrm{d}}
\newcommand{\ii}{\mathrm{i}}
\newcommand{\ee}{\mathrm{e}}
\newcommand{\lambdabar}{\bar\lambda}
\DeclareMathOperator{\tr}{tr}
\DeclareMathOperator{\rank}{rank}
\DeclareMathOperator{\Var}{Var}
\DeclareMathOperator{\Cov}{Cov}
\DeclareMathOperator{\E}{E}
\DeclareMathOperator{\Res}{Res}
\DeclareMathOperator{\sgn}{sgn}
\DeclareMathOperator{\lcm}{lcm}
\DeclareMathOperator{\ord}{ord}
\DeclareMathOperator{\Gal}{Gal}
\DeclareMathOperator{\Aut}{Aut}
\DeclareMathOperator{\diag}{diag}
\DeclareMathOperator{\adj}{adj}
\DeclareMathOperator{\erf}{erf}
\DeclareMathOperator{\erfc}{erfc}
\DeclareMathOperator{\Li}{Li}
\DeclareMathOperator{\vol}{vol}
\DeclareMathOperator{\Hom}{Hom}
\DeclareMathOperator{\Ker}{Ker}
\DeclareMathOperator{\Img}{Im}
\DeclareMathOperator*{\argmin}{arg\,min}
\DeclareMathOperator*{\argmax}{arg\,max}
\setlength{\parindent}{0pt}
\setlength{\parskip}{4pt plus 2pt}
\pagestyle{plain}
\ctexset{section={format=\Large\bfseries\color{jl},beforeskip=14pt,afterskip=8pt}}
\newcommand{\续}[1]{\InputIfFileExists{chapters/#1.tex}{}{}}
\setcounter{tocdepth}{2}
\begin{document}
{\Large\bfseries 数理积木速查：从分析代数到综合应用}\hfill{\small 2026-11-03}

{\small 每条给出条件、具体操作与结果；红框只圈触发形式、主公式或易错的符号定义。难度对标 MIT 本科高年级课程（略高于教材正文）。内容依据标准教材与课程知识整理，不是对某份资料的抽取；各节末列参考书目。其中 387 项公式与常数做了符号或数值复核，见配套《核验记录》。}

\tableofcontents

\section{分析代数}

\subsection{核心积木}

\begin{结论}{Abel 分部求和：单调因子乘有界部分和}
$B_k=\sum_{j\le k}b_j$，$|B_k|\le M$，$a_k\ge0$ 单调递减，则
\[
\sum_{k=m}^n a_kb_k=a_nB_n-a_mB_{m-1}-\sum_{k=m}^{n-1}(a_{k+1}-a_k)B_k,
\qquad \重点{\Bigl|\sum_{k=m}^n a_kb_k\Bigr|\le 2Ma_m}.
\]
配合 $\displaystyle\Bigl|\sum_{k=1}^n\sin kx\Bigr|=\Bigl|\frac{\sin\frac{nx}{2}\sin\frac{(n+1)x}{2}}{\sin\frac x2}\Bigr|\le\frac{1}{|\sin\frac x2|}$（$\cos$ 同理），得 $\sum\frac{\sin kx}{k}$ 在 $[\delta,2\pi-\delta]$ 上\textbf{一致}收敛，尾项 $\le \dfrac{2}{m\sin(\delta/2)}$。
\end{结论}

\begin{结论}{Stolz 定理求递推数列的渐近阶}
若 $x_n\to0^+$ 且 \重点{x_{n+1}=x_n-c\,x_n^{p+1}+o(x_n^{p+1})}（$c,p>0$），则 $x_{n+1}^{-p}-x_n^{-p}\to pc$，由 Stolz（$b_n=n$）
\[
\重点{x_n\sim (pcn)^{-1/p}}.
\]
例：$x_{n+1}=\sin x_n$ 有 $p=2,\ c=\tfrac16$，故 $x_n\sim\sqrt{3/n}$；$x_{n+1}=\ln(1+x_n)$ 有 $p=1,\ c=\tfrac12$，故 $x_n\sim 2/n$。
\end{结论}

\begin{结论}{Frullani 积分}
$f\in C[0,\infty)$，$f(\infty)=\lim_{x\to\infty}f(x)$ 存在，$a,b>0$：
\[
\重点{\int_0^\infty\frac{f(ax)-f(bx)}{x}\,\dd x=\bigl(f(0)-f(\infty)\bigr)\ln\frac ba}.
\]
看到“同一函数两种伸缩之差除以 $x$”即用；如 $\int_0^\infty\frac{\ee^{-ax}-\ee^{-bx}}{x}\dd x=\ln\frac ba$。
\end{结论}

\begin{结论}{插入收敛因子再对参数求导}
分母 $x$ 阻碍直接积分时，令 $I(a)=\int_0^\infty \ee^{-ax}\frac{\sin x}{x}\dd x$。对 $a\ge a_0>0$ 有控制 $\ee^{-a_0x}$，可求导：
\[
\重点{I'(a)=-\int_0^\infty \ee^{-ax}\sin x\,\dd x=-\frac{1}{1+a^2}},\qquad I(+\infty)=0\ \Rightarrow\ I(a)=\frac\pi2-\arctan a .
\]
\注意 $a\to0^+$ 的极限交换须另证：$\ee^{-ax}$ 关于 $x$ 单调且一致有界，Abel 判别给出 $[0,\infty)$ 上一致收敛，故 $\int_0^\infty\frac{\sin x}{x}\dd x=\frac\pi2$。
\end{结论}

\begin{结论}{多元 Gauss 积分（带线性项）}
$A$ 为 $n$ 阶实对称正定，$b\in\mathbb C^n$：
\[
\重点{\int_{\mathbb R^n}\exp\Bigl(-\tfrac12x^{\mathsf T}Ax+b^{\mathsf T}x\Bigr)\dd^nx=\sqrt{\frac{(2\pi)^n}{\det A}}\;\exp\Bigl(\tfrac12 b^{\mathsf T}A^{-1}b\Bigr)}.
\]
对 $b$ 求导再令 $b=0$：$\langle x_ix_j\rangle=(A^{-1})_{ij}$；高阶矩按 Wick 配对求和，$\langle x_ix_jx_kx_l\rangle=A^{-1}_{ij}A^{-1}_{kl}+A^{-1}_{ik}A^{-1}_{jl}+A^{-1}_{il}A^{-1}_{jk}$。
\end{结论}

\begin{结论}{Parseval 求 $\zeta(2k)$：注意归一化}
$f\in L^2(-\pi,\pi)$，$f\sim\frac{a_0}{2}+\sum_{n\ge1}(a_n\cos nx+b_n\sin nx)$，则
\[
\重点{\frac1\pi\int_{-\pi}^{\pi}|f|^2\dd x=\frac{a_0^2}{2}+\sum_{n\ge1}(a_n^2+b_n^2)}.
\]
$f=x$：$b_n=\frac{2(-1)^{n+1}}{n}$ 给 $\zeta(2)=\frac{\pi^2}{6}$；$f=x^2$：$a_0=\frac{2\pi^2}{3}$，$a_n=\frac{4(-1)^n}{n^2}$ 给 $\zeta(4)=\frac{\pi^4}{90}$。
\end{结论}

\begin{结论}{Dini 定理：单调逐点收敛升级为一致收敛}
$K$ \textbf{紧}，$f_n,f\in C(K)$，对每个 $x$，$f_n(x)$ \textbf{单调}趋于 $f(x)$ $\Rightarrow$ \重点{f_n\rightrightarrows f}。
三个条件缺一不可：$x^n$ 在 $[0,1]$ 上极限不连续；$K=[0,1)$ 时同样失败。常用来交换极限与积分，或证明 $\sum$ 非负连续项级数在紧集上一致收敛。
\end{结论}

\begin{结论}{秩一扰动：行列式引理与 Sherman–Morrison}
$A$ 可逆，$u,v\in\mathbb C^n$：
\[
\重点{\det(A+uv^{\mathsf T})=\det A\,(1+v^{\mathsf T}A^{-1}u)},\qquad
(A+uv^{\mathsf T})^{-1}=A^{-1}-\frac{A^{-1}uv^{\mathsf T}A^{-1}}{1+v^{\mathsf T}A^{-1}u}.
\]
取 $A=aI,\ u=b\mathbf 1,\ v=\mathbf 1$（$J=\mathbf 1\mathbf 1^{\mathsf T}$）：$aI+bJ$ 特征值为 $a$（$n-1$ 重）与 $a+nb$，$\det=a^{n-1}(a+nb)$，逆为 $\frac1a\bigl(I-\frac{b}{a+nb}J\bigr)$。
\end{结论}

\begin{结论}{分块行列式：Schur 补}
$M=\begin{pmatrix}A&B\\C&D\end{pmatrix}$，$A$ 可逆：$\重点{\det M=\det A\cdot\det(D-CA^{-1}B)}$，
\[
M^{-1}=\begin{pmatrix}A^{-1}+A^{-1}BS^{-1}CA^{-1}&-A^{-1}BS^{-1}\\-S^{-1}CA^{-1}&S^{-1}\end{pmatrix},\quad S=D-CA^{-1}B.
\]
四块同阶且 \重点{AC=CA} 时 $\det M=\det(AD-CB)$（不需 $A$ 可逆，取 $A+\varepsilon I$ 后令 $\varepsilon\to0$）。
\end{结论}

\begin{结论}{$AB$ 与 $BA$ 的特征多项式}
$A\in M_{m\times n}$，$B\in M_{n\times m}$：
\[
\重点{\lambda^{n}\det(\lambda I_m-AB)=\lambda^{m}\det(\lambda I_n-BA)},\qquad \det(I_m+AB)=\det(I_n+BA).
\]
故 $AB$、$BA$ 非零特征值连同代数重数相同。把 $n\times n$ 大矩阵的行列式降为 $m\times m$（$m\ll n$）时直接用。
\end{结论}

\begin{结论}{由秩序列读出 Jordan 结构}
对特征值 $\lambda$，令 \重点{r_k=\rank(A-\lambda I)^k}，$r_0=n$。则
\[
\#\{\text{大小}\ge k\text{ 的 }\lambda\text{ 块}\}=r_{k-1}-r_k,\qquad
\重点{\#\{\text{大小}=k\}=r_{k-1}-2r_k+r_{k+1}}.
\]
块总数 $=n-r_1=\dim\ker(A-\lambda I)$；最大块大小 $=$ 最小的 $k$ 使 $r_k=r_{k+1}$。
\end{结论}

\begin{结论}{Hermite 矩阵特征值扰动：Weyl 与交错}
特征值\textbf{降序} $\lambda_1\ge\dots\ge\lambda_n$。$A,B$ Hermite，$i+j-1\le n$：
\[
\重点{\lambda_{i+j-1}(A+B)\le\lambda_i(A)+\lambda_j(B)},\qquad |\lambda_k(A+E)-\lambda_k(A)|\le\|E\|_2 .
\]
删去一行一列得主子阵 $B$：$\lambda_k(A)\ge\lambda_k(B)\ge\lambda_{k+1}(A)$。均由 Courant–Fischer $\lambda_k=\max_{\dim V=k}\min_{x\in V}\frac{x^*Ax}{x^*x}$ 得到。
\end{结论}

\begin{结论}{Burnside 计数：项链}
有限群 $G$ 作用于 $X$：\重点{|X/G|=\frac1{|G|}\sum_{g\in G}|X^g|}。$n$ 珠 $k$ 色、只计旋转：旋转 $j$ 格的不动着色数为 $k^{\gcd(j,n)}$，
\[
\重点{N=\frac1n\sum_{d\mid n}\varphi(d)\,k^{n/d}}.
\]
再允许翻转（二面体群 $D_n$）：在旋转部分 $\sum_{d\mid n}\varphi(d)k^{n/d}$ 上，$n$ 奇时加 $n\,k^{(n+1)/2}$，$n$ 偶时加 $\frac n2(k^{n/2+1}+k^{n/2})$，再除以 $2n$。
\end{结论}

\begin{结论}{Sylow 计数与元素计数}
$|G|=p^am$，$p\nmid m$：\重点{n_p\equiv1\pmod p,\ n_p\mid m}；$n_p=1\Leftrightarrow$ Sylow $p$ 子群正规。
\begin{itemize}\setlength{\itemsep}{0pt}
\item $|G|=pq$，$p<q$，$p\nmid q-1$ $\Rightarrow$ $n_p=n_q=1$，$G\cong C_{pq}$（如 15，33，35）。
\item Sylow 子群阶为素数 $p$ 时，它们两两交平凡，贡献 \重点{n_p(p-1)} 个 $p$ 阶元。$|G|=30$：若 $n_3=10,\ n_5=6$，则已有 $20+24>30$ 个元素，矛盾，故必有正规 Sylow 子群。
\end{itemize}
\end{结论}

\begin{结论}{类方程 $\Rightarrow$ $p^2$ 阶群交换}
$|G|=|Z(G)|+\sum_i[G:C_G(x_i)]$（$x_i$ 取非中心共轭类代表）。$p$ 群中每个指标被 $p$ 整除 $\Rightarrow$ \重点{Z(G)\ne1}。
配合 \重点{G/Z(G)\text{ 循环}\Rightarrow G\text{ 交换}}：$|G|=p^2$ 时 $|Z|=p$ 会使 $G/Z$ 循环而矛盾，故 $G$ 交换，$G\cong C_{p^2}$ 或 $C_p\times C_p$。
\end{结论}

\begin{结论}{平移后用 Eisenstein：分圆多项式不可约}
$\Phi_p(x)=\frac{x^p-1}{x-1}$，$p$ 素数。代入 $x\mapsto x+1$：
\[
\重点{\Phi_p(x+1)=\sum_{k=1}^{p}\binom pk x^{k-1}},
\]
首项系数 1，其余系数被 $p$ 整除，常数项 $p$ 不被 $p^2$ 整除，Eisenstein 给出不可约，$[\mathbb Q(\zeta_p):\mathbb Q]=p-1$。一般 $[\mathbb Q(\zeta_n):\mathbb Q]=\varphi(n)$。
\end{结论}

\begin{结论}{矩阵指数的三条可用恒等式}
\[
\重点{\det\ee^{A}=\ee^{\tr A}},\qquad \ee^{A}B\,\ee^{-A}=\sum_{k\ge0}\frac{1}{k!}[A,[A,\cdots[A,B]]],
\]
若 \重点{[A,B]\text{ 与 }A,B\text{ 都对易}}，则 $\ee^A\ee^B=\ee^{A+B+\frac12[A,B]}$（量子力学中位移算符、相干态计算直接用）。
\end{结论}

\参考{Rudin《数学分析原理》；Axler, \emph{Linear Algebra Done Right}；Horn--Johnson, \emph{Matrix Analysis} 第 4 章；Artin, \emph{Algebra}（MIT 18.701）；Dummit--Foote 第 4 章。}

\续{01b}
\续{01c}
\续{01d}
\续{01e}
\续{01f}
\续{01g}
\续{01h}
\续{01i}
\续{01j}
\续{01k}
\续{01l}
\续{01m}
\续{01n}
\section{几何拓扑}

\subsection{核心积木}

\begin{结论}{Euler 示性数的四条运算}
闭曲面：可定向亏格 $g$ 时 $\chi=2-2g$，$k$ 个交叉帽时 $\chi=2-k$；连通和 $\chi(M\# N)=\chi(M)+\chi(N)-2$；$d$ 叶覆叠 \重点{\chi(\tilde X)=d\,\chi(X)}。
紧 Riemann 曲面间 $d$ 次全纯映射（Riemann--Hurwitz）：
\[
\重点{2g_X-2=d\,(2g_Y-2)+\sum_{p\in X}(e_p-1)}.
\]
例：$y^2=f(x)$，$\deg f\in\{2g+1,2g+2\}$，对 $x$ 坐标 $d=2$、分歧点 $2g+2$ 个，得亏格 $g$。
\end{结论}

\begin{结论}{粘 2-胞腔 = 加一条关系}
CW 复形 $X$ 沿闭路 $w$ 粘一个 2-胞腔：\重点{\pi_1(X\cup_w D^2)=\pi_1(X)/\langle\!\langle w\rangle\!\rangle}（van Kampen）。由此
\[
\pi_1(\Sigma_g)=\Bigl\langle a_1,b_1,\dots,a_g,b_g\ \Big|\ \prod_{i=1}^g[a_i,b_i]\Bigr\rangle,\quad H_1(\Sigma_g)=\mathbb Z^{2g},\quad \pi_1(\mathbb{RP}^2)=\mathbb Z/2 .
\]
$n$ 个圆的楔和 $\pi_1$ 为秩 $n$ 自由群；$H_1=\pi_1$ 的交换化。粘 3 维以上胞腔不改变 $\pi_1$。
\end{结论}

\begin{结论}{无收缩 $\Rightarrow$ Brouwer 不动点}
若有收缩 $r:D^2\to S^1$（$r|_{S^1}=\mathrm{id}$），则 $\mathrm{id}_{\mathbb Z}=r_*\circ i_*$ 经过 $\pi_1(D^2)=0$，矛盾。连续 $f:D^2\to D^2$ 无不动点时，\重点{\text{从 }f(x)\text{ 经 }x\text{ 的射线交 }S^1\text{ 的点}} 就是这种收缩。高维用 $H_{n-1}(S^{n-1})=\mathbb Z$ 同理。
\end{结论}

\begin{结论}{映射度：对径映射与毛球}
对径映射 $a:S^n\to S^n$ 的度 \重点{\deg a=(-1)^{n+1}}（$n+1$ 个坐标反射之积）。$f$ 无不动点 $\Rightarrow$ $f\simeq a$；$f(x)\ne -x$ 对所有 $x$ $\Rightarrow$ $f\simeq\mathrm{id}$。
若 $S^n$ 上有处处非零切向量场 $v$，则 $f_t=\cos(\pi t)x+\sin(\pi t)\frac{v(x)}{|v(x)|}$ 连接 $\mathrm{id}$ 与 $a$，需 $1=(-1)^{n+1}$：\重点{n\ \text{必须为奇数}}。
\end{结论}

\begin{结论}{Poincaré--Hopf 与 Lefschetz}
紧流形上只有孤立零点的向量场：\重点{\sum_{\text{零点}}\operatorname{ind}=\chi(M)}；平面零点指标 $=$ 绕零点小圆一周向量转过的圈数（源、汇、中心为 $+1$，鞍点为 $-1$）。
紧可三角剖分 $X$ 上，$L(f)=\sum_k(-1)^k\tr\bigl(f_*|H_k(X;\mathbb Q)\bigr)$，\重点{L(f)\ne0\Rightarrow f\text{ 有不动点}}。例：$\mathbb{RP}^{2n}$ 的有理同调只有 $H_0$，$L\equiv1$，故其上每个自映射都有不动点。
\end{结论}

\begin{结论}{提升判据：何时映射可以“展开”}
覆叠 $p:\tilde X\to X$，$Y$ 道路连通且局部道路连通，$f:Y\to X$。
\[
\重点{f\text{ 可提升}\iff f_*\pi_1(Y)\subseteq p_*\pi_1(\tilde X)}.
\]
单连通 $Y$（如 $S^n$，$n\ge2$）到 $S^1$ 的映射提升到 $\mathbb R$，故零伦；叶数 $=[\pi_1(X):p_*\pi_1(\tilde X)]$。
\end{结论}

\begin{结论}{空间曲线的曲率与挠率（任意参数）}
\[
\重点{\kappa=\frac{|\mathbf r'\times\mathbf r''|}{|\mathbf r'|^{3}},\qquad
\tau=\frac{(\mathbf r'\times\mathbf r'')\cdot\mathbf r'''}{|\mathbf r'\times\mathbf r''|^{2}}}.
\]
圆柱螺线 $(a\cos t,a\sin t,bt)$：$\kappa=\frac{a}{a^2+b^2}$，$\tau=\frac{b}{a^2+b^2}$。$\kappa,\tau$ 为常数（$\kappa>0$）的曲线只有螺线（含圆、直线退化情形）。
\end{结论}

\begin{结论}{图像曲面 $z=f(x,y)$ 的 $K$ 与 $H$}
记 $W=1+f_x^2+f_y^2$：
\[
\重点{K=\frac{f_{xx}f_{yy}-f_{xy}^2}{W^{2}}},\qquad
H=\frac{(1+f_y^2)f_{xx}-2f_xf_yf_{xy}+(1+f_x^2)f_{yy}}{2W^{3/2}}.
\]
在临界点 $\nabla f=0$ 处退化为 Hessian 的行列式与半迹。极小曲面即 $H\equiv0$。
\end{结论}

\begin{结论}{只用度量算 Gauss 曲率}
正交参数 $\dd s^2=E\,\dd u^2+G\,\dd v^2$：
\[
K=-\frac{1}{2\sqrt{EG}}\Bigl[\partial_u\Bigl(\frac{G_u}{\sqrt{EG}}\Bigr)+\partial_v\Bigl(\frac{E_v}{\sqrt{EG}}\Bigr)\Bigr].
\]
两个常用特例：
\[
\重点{\dd s^2=\lambda(\dd u^2+\dd v^2)\Rightarrow K=-\frac{1}{2\lambda}\Delta\ln\lambda},\qquad
\重点{\dd s^2=\dd r^2+G\,\dd\theta^2\Rightarrow K=-\frac{(\sqrt G)_{rr}}{\sqrt G}}.
\]
前者为共形坐标，后者为测地极坐标。
例：$\lambda=\frac{4}{(1\mp r^2)^2}$ 分别给 $K=\pm1$ 的 Poincaré 圆盘（取 $-$）与球面球极投影（取 $+$）。
\end{结论}

\begin{结论}{Gauss--Bonnet（带边、带角）}
紧定向曲面 $M$，边界分段光滑，内角 $\theta_i$：
\[
\重点{\int_M K\,\dd A+\int_{\partial M}k_g\,\dd s+\sum_i(\pi-\theta_i)=2\pi\chi(M)}.
\]
测地三角形：$\alpha+\beta+\gamma-\pi=\int_T K\,\dd A$。单位球纬线 $\theta=\theta_0$（余纬）所围极冠：$\int K=2\pi(1-\cos\theta_0)$，即该纬线上平行移动一周的转角。
\end{结论}

\begin{结论}{Clairaut 关系：旋转面测地线}
旋转面上，$\rho$ 为点到轴的距离，$\psi$ 为测地线与\textbf{纬线}夹角：\重点{\rho\cos\psi=\text{常数}}。
故测地线在 $\rho=|c|$ 处与纬线相切后折回；子午线（$\psi=\pi/2$，$c=0$）总是测地线。它是轴向角动量守恒的几何版本。
\end{结论}

\begin{结论}{闭而不恰当的形式：绕数与立体角}
$\mathbb R^2\setminus0$ 上 $\omega=\dfrac{x\,\dd y-y\,\dd x}{x^2+y^2}$ 满足 $\dd\omega=0$，但 \重点{\oint_\gamma\omega=2\pi\,\mathrm{wind}(\gamma,0)}，故非恰当。
$\mathbb R^3\setminus0$ 上 $\Omega=\dfrac{x\,\dd y\wedge\dd z+y\,\dd z\wedge\dd x+z\,\dd x\wedge\dd y}{r^3}$ 闭，$\int_{S}\Omega=4\pi$（$S$ 包围原点）。由 Stokes，任何包围原点的闭曲面积分都相同——这就是 Gauss 定律的拓扑骨架。
\end{结论}

\begin{结论}{Gauss 环绕数积分}
不相交闭曲线 $\gamma_1,\gamma_2\subset\mathbb R^3$：
\[
\重点{\mathrm{lk}(\gamma_1,\gamma_2)=\frac{1}{4\pi}\oint_{\gamma_1}\!\oint_{\gamma_2}\frac{(\mathbf r_1-\mathbf r_2)\cdot(\dd\mathbf r_1\times\dd\mathbf r_2)}{|\mathbf r_1-\mathbf r_2|^{3}}}\in\mathbb Z .
\]
物理读法：$\gamma_2$ 载单位电流，Biot--Savart 场沿 $\gamma_1$ 的环量 $=\mu_0\cdot\mathrm{lk}$（Ampère 定律）。
\end{结论}

\begin{结论}{$n$ 维球体积：用 Gauss 积分反推}
$\int_{\mathbb R^n}\ee^{-|x|^2}\dd^nx=\pi^{n/2}$，另一方面 $=|S^{n-1}|\int_0^\infty r^{n-1}\ee^{-r^2}\dd r=\tfrac12|S^{n-1}|\Gamma(\tfrac n2)$：
\[
\重点{|S^{n-1}|=\frac{2\pi^{n/2}}{\Gamma(n/2)}},\qquad V_n(R)=\frac{\pi^{n/2}}{\Gamma(\frac n2+1)}R^n .
\]
$V_n(1)$ 在 $n=5$ 最大，$n\to\infty$ 时趋于 $0$；高维球的体积集中在半径 $1-O(1/n)$ 的薄壳中。
\end{结论}

\begin{结论}{Cartan 公式：何时流保持一个形式}
\[
\重点{\mathcal L_X\omega=\dd\,\iota_X\omega+\iota_X\dd\omega}.
\]
辛形式 $\omega$（$\dd\omega=0$）与 Hamilton 场 $\iota_{X_H}\omega=\dd H$：$\mathcal L_{X_H}\omega=\dd\dd H=0$，Hamilton 流保辛；再由 $\omega^n$ 守恒得 Liouville 定理。散度版本：$\mathcal L_X(\dd V)=(\nabla\cdot X)\,\dd V$。
\end{结论}

\参考{Hatcher, \emph{Algebraic Topology} 第 1--2 章；Munkres, \emph{Topology}（MIT 18.901）；do Carmo, \emph{Differential Geometry of Curves and Surfaces}（MIT 18.950）；Spivak, \emph{Calculus on Manifolds}。}

\续{02b}
\续{02c}
\续{02d}
\续{02e}
\续{02f}
\续{02g}
\续{02h}
\section{概率论}

\subsection{核心积木}

\begin{结论}{示性变量拆期望与方差}
把计数写成 $X=\sum_i\mathbf 1_{A_i}$，不需要独立：
\[
\E X=\sum_i P(A_i),\qquad \重点{\Var X=\sum_iP(A_i)\bigl(1-P(A_i)\bigr)+\sum_{i\ne j}\bigl(P(A_iA_j)-P(A_i)P(A_j)\bigr)}.
\]
随机置换不动点数：$P(A_i)=\frac1n$，$P(A_iA_j)=\frac1{n(n-1)}$，得 $\E X=\Var X=1$；$n$ 个连续 iid 中“破纪录”次数期望 $H_n=\sum_{k\le n}\frac1k$。
\end{结论}

\begin{结论}{集卡问题：分段几何变量}
已集 $i-1$ 种后，下一新种的等待时间 $\sim\mathrm{Geom}(p_i)$，$p_i=\frac{n-i+1}{n}$，各段独立：
\[
\重点{\E T=nH_n=n\ln n+\gamma n+\tfrac12+o(1)},\qquad \Var T=\sum_{i=1}^n\frac{1-p_i}{p_i^2}<\frac{\pi^2}{6}n^2 .
\]
Chebyshev 给出 $T/(n\ln n)\to1$（依概率）。
\end{结论}

\begin{结论}{全方差公式与随机和}
\[
\重点{\Var X=\E[\Var(X\mid Y)]+\Var(\E[X\mid Y])}.
\]
$S=\sum_{i=1}^N X_i$，$X_i$ iid 且与 $N$ 独立：$\E S=\E N\cdot\E X$，\重点{\Var S=\E N\cdot\Var X+\Var N\cdot(\E X)^2}。例：$N\sim\mathrm{Poisson}(\lambda)$ 时 $\Var S=\lambda\,\E X^2$（复合 Poisson）。
\end{结论}

\begin{结论}{分支过程灭绝概率}
每个体后代数母函数 $G(s)=\sum p_ks^k$，均值 $\mu=G'(1)$，$p_1\ne1$。灭绝概率 $q$ 是
\[
\重点{s=G(s)\text{ 在 }[0,1]\text{ 上的最小根}},\qquad q=1\iff\mu\le1 .
\]
$Z_n$ 的母函数为 $G$ 的 $n$ 次复合 $G_n$；$P(Z_n=0)=G_n(0)\uparrow q$。Poisson$(\lambda)$ 后代：$q=\ee^{\lambda(q-1)}$。
\end{结论}

\begin{结论}{赌徒破产：用两条鞅直接读答案}
从 $k$ 出发，每步 $+1$ 概率 $p$、$-1$ 概率 $q=1-p$，在 $0$ 或 $N$ 停止（$\tau$）。$r=q/p\ne1$ 时 $r^{S_n}$ 与 $S_n-(p-q)n$ 是鞅，可选停时成立（$\E\tau<\infty$，增量有界）：
\[
\重点{P_k(\text{先到 }N)=\frac{1-r^{k}}{1-r^{N}}},\qquad \E_k\tau=\frac{k}{q-p}-\frac{N}{q-p}\cdot\frac{1-r^k}{1-r^N}.
\]
$p=q$ 时用 $S_n$ 与 $S_n^2-n$：概率 $k/N$，\重点{\E_k\tau=k(N-k)}。
\end{结论}

\begin{结论}{可选停时定理：哪些条件够用}
$M_n$ 鞅，$\tau$ 停时。下列之一成立则 \重点{\E M_\tau=\E M_0}：(i) $\tau\le C$；(ii) $|M_{n\wedge\tau}|\le C$；(iii) $\E\tau<\infty$ 且 $\E[|M_{n+1}-M_n|\mid\mathcal F_n]\le C$。
\注意 对称随机游走 $\tau=\inf\{n:S_n=1\}$ 几乎必然有限但 $\E S_\tau=1\ne0$，因为 $\E\tau=\infty$。
\end{结论}

\begin{结论}{Markov 链：平均返回时间 $=1/\pi_i$}
不可约正常返链：\重点{\E_i[T_i^+]=1/\pi_i}。图上简单随机游走满足细致平衡，\重点{\pi_v=\frac{\deg v}{2|E|}}。
例：国际象棋马从角格出发，棋盘马步图度数和 $336$，角格度 $2$，期望 $168$ 步回到原角。一般可逆链：$\pi_iP_{ij}=\pi_jP_{ji}$ 只需一个方程组即可求 $\pi$。
\end{结论}

\begin{结论}{集中不等式：Hoeffding 与 Chernoff}
$X_i\in[a_i,b_i]$ 独立，$S=\sum X_i$：
\[
\重点{P(S-\E S\ge t)\le\exp\Bigl(-\frac{2t^2}{\sum(b_i-a_i)^2}\Bigr)}.
\]
$X$ 为独立 Bernoulli 之和，$\mu=\E X$，$0<\delta\le1$：$P(X\ge(1+\delta)\mu)\le\ee^{-\delta^2\mu/3}$，$P(X\le(1-\delta)\mu)\le\ee^{-\delta^2\mu/2}$。推导统一为：$P(X\ge a)\le\inf_{s>0}\ee^{-sa}\E\ee^{sX}$。
\end{结论}

\begin{结论}{反射原理与选票定理}
简单随机游走，整数 $a\ge1$：\重点{P\bigl(\max_{k\le n}S_k\ge a\bigr)=2P(S_n>a)+P(S_n=a)}。Brown 运动：$P(\max_{s\le t}B_s\ge a)=2P(B_t\ge a)$，故 $\max_{s\le t}B_s\overset{d}{=}|B_t|$。
选票：$A$ 得 $a$ 票、$B$ 得 $b$ 票（$a>b$），随机计票时 $A$ 始终严格领先的概率 \重点{\frac{a-b}{a+b}}。
\end{结论}

\begin{结论}{Pólya 常返判据}
$\mathbb Z^d$ 简单随机游走：$P(S_{2n}=0)\sim c_d\,n^{-d/2}$，常返 $\iff\sum_nP(S_{2n}=0)=\infty$：
\[
\重点{d=1,2\ \text{常返};\qquad d\ge3\ \text{暂留}}.
\]
返回概率 $=1-1/\sum_{n\ge0}P(S_{2n}=0)$（$d=3$ 约 $0.34$）。
\end{结论}

\begin{结论}{指数分布的竞争与无记忆}
$X_i\sim\mathrm{Exp}(\lambda_i)$ 独立：\重点{\min_iX_i\sim\mathrm{Exp}\bigl(\sum\lambda_i\bigr)}，$P(X_j\text{ 最小})=\frac{\lambda_j}{\sum_i\lambda_i}$，且最小值与“谁最小”独立。
Poisson 过程：给定 $N(t)=n$，到达时刻与 $n$ 个 $U(0,t)$ 的次序统计量同分布；按概率 $p$ 独立着色得到速率 $p\lambda$、$(1-p)\lambda$ 的\textbf{独立} Poisson 过程。
\end{结论}

\begin{结论}{均匀次序统计量}
$U_1,\dots,U_n$ iid $U(0,1)$：
\[
\重点{U_{(k)}\sim\mathrm{Beta}(k,\,n-k+1)},\qquad \E U_{(k)}=\frac{k}{n+1},\quad \Var U_{(k)}=\frac{k(n-k+1)}{(n+1)^2(n+2)} .
\]
$n+1$ 段间隔（含两端）可交换，联合分布与 $(E_1,\dots,E_{n+1})/\sum E_j$ 相同（$E_j$ iid Exp）。
\end{结论}

\begin{结论}{二元正态：条件分布与象限概率}
$(X,Y)$ 二元正态，相关系数 $\rho$：
\[
\重点{X\mid Y=y\sim N\Bigl(\mu_X+\rho\frac{\sigma_X}{\sigma_Y}(y-\mu_Y),\ (1-\rho^2)\sigma_X^2\Bigr)}.
\]
多元：条件协方差为 Schur 补 $\Sigma_{11}-\Sigma_{12}\Sigma_{22}^{-1}\Sigma_{21}$。标准情形 \重点{P(X>0,Y>0)=\frac14+\frac{\arcsin\rho}{2\pi}}（化为极坐标下的扇形角）。
\end{结论}

\begin{结论}{Borel--Cantelli 定 $\limsup$}
$\sum P(A_n)<\infty\Rightarrow P(A_n\text{ i.o.})=0$；$A_n$ 独立且 $\sum P(A_n)=\infty\Rightarrow P(A_n\text{ i.o.})=1$。
例：$X_n$ iid $\mathrm{Exp}(1)$，取 $A_n=\{X_n>c\ln n\}$，$P(A_n)=n^{-c}$，得 \重点{\limsup_n\frac{X_n}{\ln n}=1\ \text{a.s.}}。
\end{结论}

\参考{Ross, \emph{A First Course in Probability}；Grimmett--Stirzaker, \emph{Probability and Random Processes}；Durrett, \emph{Probability: Theory and Examples}（MIT 18.600 / 18.175）；Levin--Peres, \emph{Markov Chains and Mixing Times}。}

\续{03b}
\续{03c}
\续{03d}
\续{03e}
\续{03f}
\续{03g}
\续{03h}
\续{03i}
\section{数值统计}

\subsection{核心积木}

\begin{结论}{Newton 迭代：误差递推与重根修正}
单根 $r$、$f\in C^2$、初值足够近：
\[
\重点{e_{n+1}\approx\frac{f''(r)}{2f'(r)}\,e_n^{2}}\quad(\text{二阶});\qquad \text{割线法阶 }\tfrac{1+\sqrt5}{2}.
\]
$m$ 重根时 Newton 只线性收敛，$e_{n+1}\approx\frac{m-1}{m}e_n$；改用 \重点{x_{n+1}=x_n-m\frac{f(x_n)}{f'(x_n)}}，或对 $u=f/f'$（只有单根）做 Newton，恢复二阶。
\end{结论}

\begin{结论}{复化求积误差与 Euler--Maclaurin}
步长 $h$：
\[
T_h-\int_a^bf=\frac{(b-a)h^2}{12}f''(\xi),\qquad S_h-\int_a^bf=\frac{(b-a)h^4}{180}f^{(4)}(\xi).
\]
\[
\重点{T_h-\int_a^bf=\frac{h^2}{12}\bigl[f'\bigr]_a^b-\frac{h^4}{720}\bigl[f'''\bigr]_a^b+\cdots}
\]
只含偶次幂，故 Romberg 每层消去 $h^{2j}$：$R=\frac{4^jT_{h/2}-T_h}{4^j-1}$。周期函数的端点项全消：周期 $2\pi$、在带形 $|\Im z|<a$ 内解析且 $|f|\le M$ 时，$N$ 点梯形误差 $\le\frac{4\pi M}{\ee^{aN}-1}$，\textbf{指数}收敛。
\end{结论}

\begin{结论}{Gauss 求积：$n$ 点精确到 $2n-1$ 次}
节点取 $n$ 次正交多项式（区间 $[-1,1]$、权 1 时为 Legendre $P_n$）的零点，权重 $w_i=\int\ell_i^2>0$：
\[
\重点{\int_a^bf-\sum_{i=1}^nw_if(x_i)=\frac{(b-a)^{2n+1}(n!)^4}{(2n+1)\,[(2n)!]^3}\,f^{(2n)}(\xi)}.
\]
两点：$\int_{-1}^1f\approx f(-\tfrac1{\sqrt3})+f(\tfrac1{\sqrt3})$；权 $\ee^{-x^2}$ 用 Hermite 零点，权 $(1-x^2)^{-1/2}$ 用 Chebyshev 点且权重全等于 $\pi/n$。
\end{结论}

\begin{结论}{插值误差与 Chebyshev 节点}
$n$ 个节点、次数 $\le n-1$ 的插值：$f(x)-p(x)=\dfrac{f^{(n)}(\xi)}{n!}\prod_{k=1}^n(x-x_k)$。在 $[-1,1]$ 取 \重点{x_k=\cos\frac{(2k-1)\pi}{2n}}，则 $\prod(x-x_k)=2^{1-n}T_n(x)$，
\[
\重点{\max_{[-1,1]}\Bigl|\prod_{k=1}^n(x-x_k)\Bigr|=2^{1-n}},
\]
是所有首一 $n$ 次多项式中最小的。等距节点对 $\frac{1}{1+25x^2}$ 发散（Runge 现象）。
\end{结论}

\begin{结论}{条件数：误差放大与“不要解正规方程”}
$Ax=b$：\重点{\frac{\|\delta x\|}{\|x\|}\le\kappa(A)\frac{\|\delta b\|}{\|b\|}}，$\kappa_2=\sigma_{\max}/\sigma_{\min}$。后向稳定算法的前向误差 $\approx\kappa\cdot\varepsilon_{\mathrm{mach}}$。
最小二乘：正规方程 $A^{\mathsf T}Ax=A^{\mathsf T}b$ 的条件数为 \重点{\kappa(A)^2}；QR 或 SVD 只受 $\kappa(A)$（加残差项）影响。
\end{结论}

\begin{结论}{共轭梯度的收敛率}
$A$ 对称正定，$\kappa=\lambda_{\max}/\lambda_{\min}$：
\[
\重点{\|e_k\|_A\le2\Bigl(\frac{\sqrt\kappa-1}{\sqrt\kappa+1}\Bigr)^{k}\|e_0\|_A}.
\]
达到精度 $\varepsilon$ 约需 $\frac{\sqrt\kappa}{2}\ln\frac2\varepsilon$ 步（最速下降需 $\sim\frac\kappa2\ln\frac1\varepsilon$）。$A$ 只有 $m$ 个不同特征值时至多 $m$ 步精确。预条件的目标是让 $M^{-1}A$ 的谱聚集。
\end{结论}

\begin{结论}{幂法族的收敛速率}
幂法误差 $\sim|\lambda_2/\lambda_1|^k$；位移逆迭代用 $(A-\sigma I)^{-1}$，速率 \重点{\frac{|\lambda_j-\sigma|}{|\lambda_k-\sigma|}}（$\lambda_j$ 为离 $\sigma$ 最近者，$\lambda_k$ 次近）。
Rayleigh 商迭代 $\sigma_k=\frac{x_k^{*}Ax_k}{x_k^{*}x_k}$：对 Hermite 矩阵\textbf{三阶}收敛；Rayleigh 商本身误差是特征向量误差的平方。
\end{结论}

\begin{结论}{ODE 方法的绝对稳定性}
试验方程 $y'=\lambda y$，$z=h\lambda$：
\[
\text{前向 Euler: } \重点{|1+z|\le1};\quad \text{后向 Euler: } \Bigl|\frac{1}{1-z}\Bigr|\le1\ (\text{A-稳定});\quad \text{RK4: } \Bigl|\sum_{k=0}^4\frac{z^k}{k!}\Bigr|\le1 .
\]
$y'=-ky$ 用前向 Euler 需 $h<2/k$；RK4 实轴稳定区间约 $[-2.785,0]$。刚性问题（$|\lambda|$ 差几个量级）中步长受最快模式限制，应改用隐式方法。
\end{结论}

\begin{结论}{von Neumann 分析：热方程与波方程差分}
代入 $u_j^n=G^n\ee^{\ii jk\Delta x}$。显式 FTCS 热方程，$r=D\Delta t/\Delta x^2$：
\[
\重点{G=1-4r\sin^2\tfrac{k\Delta x}{2},\qquad |G|\le1\iff r\le\tfrac12}.
\]
Crank--Nicolson 无条件稳定，二阶；波方程中心差分需 CFL 条件 \重点{c\Delta t/\Delta x\le1}；一阶迎风格式对 $u_t+cu_x=0$（$c>0$）稳定条件相同。
\end{结论}

\begin{结论}{正态样本：$\bar X$ 与 $S^2$ 的分布}
$X_i$ iid $N(\mu,\sigma^2)$，$S^2=\frac1{n-1}\sum(X_i-\bar X)^2$：
\[
\重点{\bar X\perp S^2},\qquad \frac{(n-1)S^2}{\sigma^2}\sim\chi^2_{n-1},\qquad \frac{\bar X-\mu}{S/\sqrt n}\sim t_{n-1}.
\]
$\E S^2=\sigma^2$ 对任意有方差的分布都成立，但独立性与 $\chi^2$ 分布\textbf{只对正态}成立；$\Var S^2=\frac{2\sigma^4}{n-1}$（正态）。
\end{结论}

\begin{结论}{MLE 的渐近正态与 Cramér--Rao}
Fisher 信息 $I(\theta)=\E[(\partial_\theta\ln f)^2]=-\E[\partial_\theta^2\ln f]$（正则条件下）：
\[
\重点{\sqrt n(\hat\theta_{\mathrm{MLE}}-\theta)\Rightarrow N\bigl(0,I(\theta)^{-1}\bigr)},\qquad \Var\hat\theta\ge\frac{1}{nI(\theta)}\ (\text{无偏}).
\]
Poisson$(\lambda)$：$I=1/\lambda$；Exp（速率 $\lambda$）：$I=1/\lambda^2$；$N(\mu,\sigma^2)$ 对 $\mu$：$I=1/\sigma^2$。
\注意 $U(0,\theta)$ 支撑依赖参数，不满足正则条件，MLE 为 $X_{(n)}$，误差阶 $1/n$ 而非 $1/\sqrt n$。
\end{结论}

\begin{结论}{Delta 方法与方差稳定}
$\sqrt n(T_n-\mu)\Rightarrow N(0,\sigma^2)$，$g'(\mu)\ne0$：\重点{\sqrt n\bigl(g(T_n)-g(\mu)\bigr)\Rightarrow N\bigl(0,g'(\mu)^2\sigma^2\bigr)}。
选 $g'\propto1/\sigma(\mu)$ 使方差与参数无关：Poisson 计数 $\sqrt X$ 方差 $\approx\frac14$；比例 $\arcsin\sqrt{\hat p}$ 方差 $\approx\frac1{4n}$；对数几率 $\ln\frac{\hat p}{1-\hat p}$ 方差 $\approx\frac{1}{np(1-p)}$。
\end{结论}

\begin{结论}{线性回归的协方差与自由度}
$y=X\beta+\varepsilon$，$X$ 为 $n\times p$ 满列秩，$\varepsilon\sim N(0,\sigma^2I)$：
\[
\hat\beta=(X^{\mathsf T}X)^{-1}X^{\mathsf T}y,\quad \重点{\Cov\hat\beta=\sigma^2(X^{\mathsf T}X)^{-1}},\quad \hat\sigma^2=\frac{\mathrm{RSS}}{n-p}.
\]
帽子矩阵 $H=X(X^{\mathsf T}X)^{-1}X^{\mathsf T}$ 有 $\tr H=p$，这就是减 $p$ 的来源。一元情形 $\mathrm{SE}(\hat\beta_1)=\hat\sigma/\sqrt{\sum(x_i-\bar x)^2}$：拉开 $x$ 的取值范围比加样本更省。
\end{结论}

\begin{结论}{似然比检验与 Pearson $\chi^2$ 的自由度}
Wilks：$H_0$ 下 \重点{-2\ln\Lambda\Rightarrow\chi^2_{d}}，$d=\dim\Theta-\dim\Theta_0$（真值在内点）。
Pearson 拟合优度 $\sum\frac{(O_i-E_i)^2}{E_i}$，自由度 \重点{k-1-(\text{估计参数个数})}；$r\times c$ 列联表独立性检验自由度 $(r-1)(c-1)$。期望频数小于约 5 的格应合并。
\end{结论}

\begin{结论}{共轭先验：后验只是更新参数}
Beta$(\alpha,\beta)$ 先验 $+$ $n$ 次试验 $k$ 次成功 $\Rightarrow$ \重点{\mathrm{Beta}(\alpha+k,\ \beta+n-k)}。
正态均值：先验 $N(\mu_0,\tau^2)$，$n$ 个方差 $\sigma^2$ 已知的观测：后验\textbf{精度相加} $\frac1{\tau_n^2}=\frac1{\tau^2}+\frac n{\sigma^2}$，均值 $\mu_n=\tau_n^2\bigl(\frac{\mu_0}{\tau^2}+\frac{n\bar x}{\sigma^2}\bigr)$。Gamma 先验配 Poisson：$\mathrm{Gamma}(\alpha,\beta)\to\mathrm{Gamma}(\alpha+\sum x_i,\beta+n)$（$\beta$ 为速率）。
\end{结论}

\参考{Trefethen--Bau, \emph{Numerical Linear Algebra}；Süli--Mayers, \emph{An Introduction to Numerical Analysis}（MIT 18.330 / 18.335）；Trefethen--Weideman (2014) \emph{SIAM Review} 梯形法则综述；Casella--Berger, \emph{Statistical Inference}；Wasserman, \emph{All of Statistics}（MIT 18.650）。}

\续{04b}
\续{04c}
\续{04d}
\续{04e}
\续{04f}
\续{04g}
\续{04h}
\续{04i}
\section{应用数学}

\subsection{核心积木}

\begin{结论}{Laplace 方法与 Stirling}
$f$ 在内点 $x_0$ 取唯一最大值，$f''(x_0)<0$，$g(x_0)\ne0$，$\lambda\to\infty$：
\[
\重点{\int_a^b g(x)\ee^{\lambda f(x)}\dd x\sim g(x_0)\,\ee^{\lambda f(x_0)}\sqrt{\frac{2\pi}{\lambda|f''(x_0)|}}}.
\]
$n!=n^{n+1}\int_0^\infty\ee^{n(\ln t-t)}\dd t$，$t_0=1$，$f''=-1$：$n!\sim\sqrt{2\pi n}\,(n/e)^n$。极值在端点且 $f'(x_0)\ne0$ 时改为 $\frac{g(x_0)\ee^{\lambda f(x_0)}}{\lambda|f'(x_0)|}$。
\end{结论}

\begin{结论}{驻相法}
$f'(x_0)=0$，$f''(x_0)\ne0$，其余处 $f'\ne0$：
\[
\重点{\int g(x)\ee^{\ii\lambda f(x)}\dd x\sim g(x_0)\sqrt{\frac{2\pi}{\lambda|f''(x_0)|}}\;\ee^{\ii\lambda f(x_0)+\ii\frac\pi4\operatorname{sgn}f''(x_0)}}.
\]
无驻点时由分部积分只剩端点贡献 $O(1/\lambda)$。物理中给出群速度：波包 $\int A(k)\ee^{\ii(kx-\omega t)}$ 在 $x/t=\omega'(k)$ 处振幅 $\propto t^{-1/2}$。
\end{结论}

\begin{结论}{留数算实积分的两种闭合方式}
有理函数 $R$，分母次数至少高 1，$k>0$（Jordan 引理）：
\[
\重点{\int_{-\infty}^{\infty}R(x)\ee^{\ii kx}\dd x=2\pi\ii\sum_{\Im z_j>0}\Res_{z_j}\bigl(R(z)\ee^{\ii kz}\bigr)},\qquad \int_{-\infty}^{\infty}\frac{\cos kx}{x^2+a^2}\dd x=\frac{\pi}{a}\ee^{-ka}.
\]
三角有理式：$z=\ee^{\ii\theta}$，$\cos\theta=\frac{z+z^{-1}}2$，$\dd\theta=\frac{\dd z}{\ii z}$，得 $\int_0^{2\pi}\frac{\dd\theta}{a+b\cos\theta}=\frac{2\pi}{\sqrt{a^2-b^2}}$（$a>|b|$）。
\end{结论}

\begin{结论}{用 $\pi\cot\pi z$ 求级数和}
$f$ 有理、在整数点无极点、$|zf(z)|\to0$：
\[
\重点{\sum_{n\in\mathbb Z}f(n)=-\sum_{f\text{ 的极点 }z_j}\Res_{z_j}\bigl(\pi\cot(\pi z)\,f(z)\bigr)}.
\]
例：$\sum_{n\in\mathbb Z}\frac{1}{n^2+a^2}=\frac\pi a\coth\pi a$；令 $a\to0$ 并减去 $n=0$ 项再得 $\zeta(2)$。交错和把 $\cot$ 换成 $\csc$。
\end{结论}

\begin{结论}{两点边值问题的 Green 函数}
$Lu=-(pu')'+qu=f$，$u_1$ 满足左边界条件，$u_2$ 满足右边界条件，$p\,(u_1'u_2-u_1u_2')\equiv C\ne0$：
\[
\重点{G(x,s)=\frac{u_1(x_<)\,u_2(x_>)}{C}},\qquad u(x)=\int G(x,s)f(s)\,\dd s .
\]
$-u''=f$，$u(0)=u(1)=0$：$G=x_<(1-x_>)$。若齐次问题有非零解（$C=0$），须先满足可解性条件 $\int fu_0=0$。
\end{结论}

\begin{结论}{热方程：分离变量与最慢模}
$u_t=Du_{xx}$，$u(0,t)=u(L,t)=0$：
\[
u=\sum_{n\ge1}b_n\ee^{-D(n\pi/L)^2t}\sin\frac{n\pi x}{L},\qquad b_n=\frac2L\int_0^Lu_0\sin\frac{n\pi x}{L}\dd x .
\]
长时间由第一模主导，弛豫时间 \重点{\tau=\frac{L^2}{\pi^2D}}；Neumann 边界换成 $\cos$ 级数，均值 $\frac1L\int u_0$ 守恒。
\end{结论}

\begin{结论}{热核与镜像法}
全直线：$u(x,t)=\int\dfrac{\ee^{-(x-y)^2/4Dt}}{\sqrt{4\pi Dt}}u_0(y)\dd y$。半直线 $x>0$：
\[
\重点{\text{Dirichlet }u(0,t)=0:\ \text{奇延拓}};\qquad \text{Neumann }u_x(0,t)=0:\ \text{偶延拓}.
\]
源项用 Duhamel：$u=\int_0^t\!\int K(x-y,t-s)F(y,s)\dd y\,\dd s$。扩散距离 $\sim\sqrt{2Dt}$（每维方差 $2Dt$）。
\end{结论}

\begin{结论}{d'Alembert 公式与依赖域}
$u_{tt}=c^2u_{xx}$，$u(x,0)=f$，$u_t(x,0)=g$：
\[
\重点{u(x,t)=\frac{f(x-ct)+f(x+ct)}{2}+\frac1{2c}\int_{x-ct}^{x+ct}g(s)\,\dd s}.
\]
半无界弦固定端 $u(0,t)=0$：奇延拓（反射波反相）；自由端偶延拓。三维的 Kirchhoff 公式只依赖球面 $|y-x|=ct$（Huygens 原理），二维则依赖整个圆盘（有尾迹）。
\end{结论}

\begin{结论}{Laplace 方程：Poisson 核}
圆盘 $r<R$，边值 $f(\phi)$：
\[
\重点{u(r,\theta)=\frac{1}{2\pi}\int_0^{2\pi}\frac{R^2-r^2}{R^2-2Rr\cos(\theta-\phi)+r^2}f(\phi)\,\dd\phi}.
\]
$r=0$ 给均值性质；由此得最大值原理与 Dirichlet 问题解的唯一性。等价的级数形式 $u=\frac{a_0}{2}+\sum(r/R)^n(a_n\cos n\theta+b_n\sin n\theta)$。
\end{结论}

\begin{结论}{量纲分析：$\Pi$ 定理}
$n$ 个物理量、量纲矩阵秩 $r$ $\Rightarrow$ 关系只依赖 \重点{n-r} 个无量纲组。
点爆炸（G.\,I.\,Taylor）：$R,E,t,\rho$，秩 3，唯一无量纲组 $\frac{R^5\rho}{Et^2}$，故 \重点{R=C\,(Et^2/\rho)^{1/5}}。单摆：$T=\sqrt{\ell/g}\,\Phi(\theta_0)$，量纲分析不能确定 $\Phi$ 但能排除质量。
\end{结论}

\begin{结论}{边界层：匹配渐近展开}
$\varepsilon y''+y'+y=0$，$y(0)=0$，$y(1)=1$。$y'$ 系数为正 $\Rightarrow$ 层在\textbf{左端}，内变量 $X=x/\varepsilon$：
\[
y_{\mathrm{out}}=\ee^{1-x},\quad y_{\mathrm{in}}=\ee\,(1-\ee^{-X}),\quad
\重点{y\approx \ee^{1-x}-\ee^{1-x/\varepsilon}}.
\]
匹配条件 $\lim_{X\to\infty}y_{\mathrm{in}}=\lim_{x\to0}y_{\mathrm{out}}$；复合解 $=$ 外解 $+$ 内解 $-$ 公共部分。层厚 $\sim\varepsilon$；若方程为 $\varepsilon y''-y'+\dots$，层移到右端。
\end{结论}

\begin{结论}{平面不动点：迹--行列式分类}
Jacobi 矩阵 $J$，$\tau=\tr J$，$\Delta=\det J$：
\[
\重点{\Delta<0\Rightarrow\text{鞍点}};\quad \Delta>0:\ \tau^2<4\Delta\ \text{焦点},\ \tau^2>4\Delta\ \text{结点};\quad \text{稳定}\iff\tau<0<\Delta .
\]
$\tau=0<\Delta$ 时线性化只给中心，非线性项决定实际行为。参数穿过 $\tau=0$（$\Delta>0$）时一对特征值穿过虚轴：Hopf 分岔，产生极限环。
\end{结论}

\begin{结论}{变分法：Beltrami 首次积分}
$J=\int F(y,y')\dd x$ 中 \重点{F\text{ 不显含 }x} 时，Euler--Lagrange 方程化为
\[
\重点{F-y'F_{y'}=C}.
\]
最速降线 $F=\sqrt{(1+y'^2)/y}$（$y$ 向下）：$y(1+y'^2)=2a$，即摆线 $x=a(\theta-\sin\theta)$，$y=a(1-\cos\theta)$。约束 $\int G=\ell$ 时对 $F+\lambda G$ 用同一公式。
\end{结论}

\begin{结论}{特征线、激波与 Rankine--Hugoniot}
$u_t+f(u)_x=0$：沿 $\dot x=f'(u)$ 时 $u$ 为常数。Burgers（$f=\frac{u^2}{2}$）的破碎时间 \重点{t_b=-\frac{1}{\min_xu_0'(x)}}。
激波速度 \重点{s=\frac{f(u_L)-f(u_R)}{u_L-u_R}}，Burgers 为 $\frac{u_L+u_R}2$；熵条件 $f'(u_L)>s>f'(u_R)$，否则形成稀疏波 $u=x/t$。
\end{结论}

\参考{Bender--Orszag, \emph{Advanced Mathematical Methods for Scientists and Engineers}（MIT 18.305）；Strang, \emph{Computational Science and Engineering}（18.085）；Strauss, \emph{Partial Differential Equations}；Strogatz, \emph{Nonlinear Dynamics and Chaos}；Ablowitz--Fokas, \emph{Complex Variables}。}

\续{05b}
\续{05c}
\续{05d}
\续{05e}
\续{05f}
\续{05g}
\续{05h}
\section{力学}

\subsection{核心积木}

\begin{结论}{变质量：火箭方程}
喷气相对火箭速度 $u$（大小），外力 $F$：\重点{m\dot v=-u\,\dot m+F}（$\dot m<0$）。
\[
F=0:\ \重点{\Delta v=u\ln\frac{m_0}{m_f}};\qquad \text{竖直重力场：}\ v(t)=u\ln\frac{m_0}{m(t)}-gt .
\]
\注意 不能对系统直接写 $F=\frac{\dd(mv)}{\dd t}$：喷出的质量带走动量 $\dd m\,(v-u)$，必须对“火箭 $+$ 喷出物”取动量差。
\end{结论}

\begin{结论}{有效势与圆轨道稳定性}
有心力，单位质量角动量 $h=r^2\dot\theta$：\重点{U_{\mathrm{eff}}(r)=U(r)+\frac{h^2}{2r^2}}。幂律力 $F=-kr^{\alpha}$（$k>0$）：
\[
\text{圆轨道稳定}\iff\alpha>-3,\qquad \重点{\frac{\omega_r}{\omega_\theta}=\sqrt{3+\alpha}}.
\]
近日点进动角每周 $2\pi(1/\sqrt{3+\alpha}-1)$；所有有界轨道闭合只发生在 $\alpha=1$（$\omega_r=2\omega_\theta$）和 $\alpha=-2$（$\omega_r=\omega_\theta$）——Bertrand 定理。
\end{结论}

\begin{结论}{轨道方程（Binet）与圆锥曲线参数}
$u=1/r$：\重点{u''+u=-\frac{F(1/u)}{h^2u^2}}（$F$ 为单位质量径向力）。平方反比 $F=-GM/r^2$：
\[
r=\frac{p}{1+e\cos\theta},\qquad p=\frac{h^2}{GM},\qquad \重点{e^2=1+\frac{2Eh^2}{(GM)^2}}\ (E\text{ 为单位质量能量}).
\]
由此反推力律：已知轨道 $r(\theta)$，代入左端即得 $F$。
\end{结论}

\begin{结论}{Kepler 问题的三条速算式}
\[
\重点{v^2=GM\Bigl(\frac2r-\frac1a\Bigr)},\qquad E=-\frac{GMm}{2a},\qquad T^2=\frac{4\pi^2a^3}{G(M+m)} .
\]
Hohmann 转移（圆轨道 $r_1\to r_2$）：$\Delta v_1=\sqrt{\frac{GM}{r_1}}\Bigl(\sqrt{\frac{2r_2}{r_1+r_2}}-1\Bigr)$，$\Delta v_2=\sqrt{\frac{GM}{r_2}}\Bigl(1-\sqrt{\frac{2r_1}{r_1+r_2}}\Bigr)$，转移时间为半个椭圆周期 $\pi\sqrt{(r_1+r_2)^3/(8GM)}$。
\end{结论}

\begin{结论}{纯滚动：加速度与所需摩擦}
记 \重点{\beta=I_{\mathrm{cm}}/(mR^2)}（实心球 $\frac25$，实心柱 $\frac12$，薄环 $1$）。斜面倾角 $\theta$ 纯滚：
\[
a=\frac{g\sin\theta}{1+\beta},\qquad f=\frac{\beta}{1+\beta}mg\sin\theta,\qquad \重点{\mu_s\ge\frac{\beta\tan\theta}{1+\beta}}.
\]
对\textbf{接触点}取力矩可一步得 $a$：$mgR\sin\theta=(I_{\mathrm{cm}}+mR^2)\frac aR$（只在接触点加速度指向质心时成立；均匀圆体纯滚满足此条件）。
\end{结论}

\begin{结论}{滑动到纯滚：对接触点角动量守恒}
摩擦力过地面接触线，重力与支持力力矩相消，故对地面上固定点的角动量守恒：
\[
\重点{v_f=\frac{v_0+\beta R\omega_0}{1+\beta}}\qquad(\omega_0\text{ 与前进方向同向为正}).
\]
保龄球（实心球、$\omega_0=0$）：$v_f=\frac57v_0$，滑动时间 $t=\frac{2v_0}{7\mu g}$，损失动能 $\frac27\cdot\frac12mv_0^2$。
\end{结论}

\begin{结论}{二体碰撞：质心系与散射角上限}
一维弹性：$v_1'=\dfrac{(m_1-m_2)v_1+2m_2v_2}{m_1+m_2}$。动能 $=\frac12MV^2+\frac12\mu v_{\mathrm{rel}}^2$，$\mu=\frac{m_1m_2}{m_1+m_2}$；
恢复系数 $e$ 时损失 \重点{\Delta K=\frac12\mu(1-e^2)v_{\mathrm{rel}}^2}。弹性碰撞中 $m_1>m_2$ 撞静止 $m_2$，入射粒子偏转角 \重点{\sin\theta_{\max}=m_2/m_1}；等质量时两出射方向垂直。
\end{结论}

\begin{结论}{转动参考系：四项惯性力}
\[
\重点{m\ddot{\mathbf r}=\mathbf F-2m\boldsymbol\Omega\times\dot{\mathbf r}-m\boldsymbol\Omega\times(\boldsymbol\Omega\times\mathbf r)-m\dot{\boldsymbol\Omega}\times\mathbf r}.
\]
纬度 $\lambda$ 处自高 $h$ 静止下落，向东偏 $\frac13\Omega g t^3\cos\lambda=\frac13\Omega\cos\lambda\sqrt{8h^3/g}$；Foucault 摆进动角速度 $\Omega\sin\lambda$。
\end{结论}

\begin{结论}{广义坐标下的小振动频率}
$L=\frac12M(q)\dot q^2-U_{\mathrm{eff}}(q)$，$U_{\mathrm{eff}}'(q_0)=0$：\重点{\omega^2=\frac{U_{\mathrm{eff}}''(q_0)}{M(q_0)}}。
例：半径 $R$ 的圆环绕竖直直径以 $\Omega$ 转动，环上珠子（$\theta$ 自最低点量起），$U_{\mathrm{eff}}=-\frac12mR^2\Omega^2\sin^2\theta-mgR\cos\theta$。$\Omega^2<g/R$ 时底部稳定，$\omega^2=g/R-\Omega^2$；$\Omega^2>g/R$ 时稳定点 $\cos\theta_0=\frac{g}{\Omega^2R}$，\重点{\omega^2=\Omega^2\sin^2\theta_0}。
\end{结论}

\begin{结论}{受迫阻尼振子：振幅、相位与 $Q$}
$\ddot x+\gamma\dot x+\omega_0^2x=\frac{F_0}{m}\cos\omega t$：
\[
\重点{A=\frac{F_0/m}{\sqrt{(\omega_0^2-\omega^2)^2+\gamma^2\omega^2}}},\qquad \tan\delta=\frac{\gamma\omega}{\omega_0^2-\omega^2}.
\]
$Q=\omega_0/\gamma$；吸收功率曲线半高全宽 $=\gamma$；振幅峰在 $\sqrt{\omega_0^2-\gamma^2/2}$（需 $Q>1/\sqrt2$）；自由振动能量按 $\ee^{-\gamma t}$ 衰减，每周期损失比例 $2\pi/Q$。
\end{结论}

\begin{结论}{位力定理}
有界运动、$U\propto r^n$：\重点{2\langle T\rangle=n\langle U\rangle}。引力（$n=-1$）：$E=-\langle T\rangle=\frac12\langle U\rangle$。
引力束缚系统失去能量会\textbf{升温}（负热容）；星系团速度弥散 $\sigma$、半径 $R$ 给出质量 $M\sim\sigma^2R/G$。谐振子（$n=2$）：$\langle T\rangle=\langle U\rangle$。
\end{结论}

\begin{结论}{潮汐加速度与 Hill 半径}
质量 $M$ 在距离 $r$ 处，对尺度 $d\ll r$ 的物体：沿连线潮汐加速度 \重点{\frac{2GMd}{r^{3}}}，垂直方向 $-\frac{GMd}{r^3}$。
小天体 $m$ 绕 $M$ 以半径 $a$ 公转，引力主导区域：\重点{r_H\approx a\Bigl(\frac{m}{3M}\Bigr)^{1/3}}。刚体卫星的 Roche 极限 $d\approx R\,(2\rho_M/\rho_m)^{1/3}$（流体约 $2.44R(\rho_M/\rho_m)^{1/3}$）。
\end{结论}

\begin{结论}{陀螺进动与复摆}
快速陀螺（$I_3\omega_3\gg$ 其他角动量）：\重点{\Omega_p=\frac{Mg\ell}{I_3\omega_3}}，方向由 $\boldsymbol\tau=\boldsymbol\Omega_p\times\mathbf L$ 定。
复摆：$T=2\pi\sqrt{\frac{I_P}{mgd}}=2\pi\sqrt{\frac{k^2+d^2}{gd}}$（$k$ 为回转半径），在 \重点{d=k} 时最短；等效摆长 $d+k^2/d$ 处为打击中心。大振幅修正 $T\approx T_0\bigl(1+\frac{\theta_0^2}{16}\bigr)$。
\end{结论}

\begin{结论}{绳索：绞盘公式与悬链线}
绕柱包角 $\theta$、摩擦系数 $\mu$，将滑未滑：\重点{T_{\text{大}}=T_{\text{小}}\,\ee^{\mu\theta}}（只与包角有关，与柱半径无关）。
均匀重链，最低点张力 $T_0$，线密度 $\lambda$：$y=a\cosh\frac xa$，\重点{a=\frac{T_0}{\lambda g}}，任一点张力 $T=\lambda g\,y$。
\end{结论}

\参考{Kleppner--Kolenkow, \emph{An Introduction to Mechanics}（MIT 8.012）；Morin, \emph{Introduction to Classical Mechanics}；Taylor, \emph{Classical Mechanics}；Murray--Dermott, \emph{Solar System Dynamics}。}

\续{06b}
\续{06c}
\续{06d}
\续{06e}
\section{四大力学}

\subsection{理论力学}

\begin{结论}{带电粒子的 Lagrange 量与正则动量}
\[
\重点{L=\tfrac12m\mathbf v^2-q\phi+q\,\mathbf v\cdot\mathbf A},\qquad \mathbf p=m\mathbf v+q\mathbf A,\qquad H=\frac{(\mathbf p-q\mathbf A)^2}{2m}+q\phi .
\]
守恒的是\textbf{正则}动量：均匀 $B\hat z$ 中取 $\mathbf A=Bx\hat y$，$y$ 为循环坐标，$p_y=m\dot y+qBx$ 守恒，直接给出回旋中心 $x_0=p_y/(qB)$。
\end{结论}

\begin{结论}{简正模：$\det(K-\omega^2M)=0$}
$L=\frac12\dot q^{\mathsf T}M\dot q-\frac12q^{\mathsf T}Kq$：\重点{\det(K-\omega^2M)=0}，模态对 $M$ 正交。
两单摆（摆长 $\ell$）由弹簧 $k$ 相连：$\omega_1=\sqrt{g/\ell}$（同相），$\omega_2=\sqrt{g/\ell+2k/m}$（反相）；从单侧起振时能量以拍频 $\omega_2-\omega_1$ 来回转移。$N$ 质点链（间距 $a$）：\重点{\omega(q)=2\sqrt{k/m}\,\bigl|\sin\tfrac{qa}{2}\bigr|}。
\end{结论}

\begin{结论}{绝热不变量 $J=\oint p\,\dd q$}
参数变化时间尺度远大于周期时 \重点{J=\oint p\,\dd q} 近似守恒；谐振子 $J=2\pi E/\omega$，即 $E/\omega$ 不变。
缓慢缩短摆长：$E\propto\ell^{-1/2}$，振幅 $\theta_0\propto\ell^{-3/4}$。磁场中回旋：磁矩 $\mu=\frac{mv_\perp^2}{2B}$ 不变，由此得磁镜反射条件 $\sin^2\alpha_0>B_0/B_{\max}$。
\end{结论}

\begin{结论}{Euler 方程：中间轴不稳定}
主轴系：$I_1\dot\omega_1=(I_2-I_3)\omega_2\omega_3$（轮换）。绕轴 2 以 $\Omega$ 转动时线性化：
\[
\ddot\omega_1=\重点{\frac{(I_1-I_2)(I_2-I_3)}{I_1I_3}}\Omega^2\omega_1 .
\]
$I_2$ 居中时系数为正，扰动指数增长（网球拍定理）。对称陀螺 $I_1=I_2$ 自由转动：体坐标中 $\boldsymbol\omega$ 绕对称轴以 $\frac{I_3-I_1}{I_1}\omega_3$ 进动。
\end{结论}

\subsection{电动力学}

\begin{结论}{辐射功率：Larmor 与 Liénard}
非相对论：\重点{P=\frac{q^2a^2}{6\pi\varepsilon_0c^3}}。相对论（实验室系）：
\[
P=\frac{q^2\gamma^6}{6\pi\varepsilon_0c^3}\Bigl(a^2-\frac{|\mathbf v\times\mathbf a|^2}{c^2}\Bigr).
\]
$\mathbf a\parallel\mathbf v$ 时 $\propto\gamma^6a^2$；圆周运动 $\mathbf a\perp\mathbf v$ 时 $\propto\gamma^4a^2$，同步辐射每圈能损 $\propto\gamma^4/R$，故电子环比质子环受限严重 $(m_p/m_e)^4$ 倍。
\end{结论}

\begin{结论}{偶极辐射与短天线辐射电阻}
$p(t)=p_0\cos\omega t$，远场 $\frac{\dd P}{\dd\Omega}\propto\sin^2\theta$：
\[
\重点{\langle P\rangle=\frac{\mu_0p_0^2\omega^4}{12\pi c}},\qquad \langle P\rangle_{\text{磁偶极}}=\frac{\mu_0m_0^2\omega^4}{12\pi c^3}.
\]
长 $\ell\ll\lambda$、电流均匀 $I_0$ 的短天线：$p_0=I_0\ell/\omega$，\重点{R_{\mathrm{rad}}=\frac{2\pi}{3}Z_0\Bigl(\frac\ell\lambda\Bigr)^2\approx790\,(\ell/\lambda)^2\ \Omega}。$\omega^4$ 依赖解释天空蓝色（Rayleigh 散射）。
\end{结论}

\begin{结论}{Maxwell 应力张量求力}
\[
\重点{T_{ij}=\varepsilon_0\bigl(E_iE_j-\tfrac12\delta_{ij}E^2\bigr)+\tfrac1{\mu_0}\bigl(B_iB_j-\tfrac12\delta_{ij}B^2\bigr)},\qquad
\mathbf F=\oint_ST\cdot\dd\mathbf a-\varepsilon_0\mu_0\frac{\dd}{\dd t}\int_V\mathbf S\,\dd\tau .
\]
静场时取一个把物体包住、场易算的曲面（如对称面加无穷远半球面）。导体表面受力面密度 $\frac{\sigma^2}{2\varepsilon_0}$；平面波辐射压：全吸收 $I/c$，全反射 $2I/c$。
\end{结论}

\begin{结论}{场的 Lorentz 变换与不变量}
参考系以 $\mathbf v$ 运动：
\[
\mathbf E'_\parallel=\mathbf E_\parallel,\quad \重点{\mathbf E'_\perp=\gamma(\mathbf E+\mathbf v\times\mathbf B)_\perp},\quad \mathbf B'_\parallel=\mathbf B_\parallel,\quad \mathbf B'_\perp=\gamma\Bigl(\mathbf B-\frac{\mathbf v\times\mathbf E}{c^2}\Bigr)_\perp .
\]
不变量 \重点{E^2-c^2B^2,\ \mathbf E\cdot\mathbf B}。若 $\mathbf E\cdot\mathbf B=0$ 且 $E<cB$，存在一个参考系只有磁场（速度 $\frac{\mathbf E\times\mathbf B}{B^2}$），即 $\mathbf E\times\mathbf B$ 漂移。
\end{结论}

\begin{结论}{匀速运动电荷的场：压扁的 Coulomb 场}
电荷以 $v$ 匀速运动，$\theta$ 为场点相对\textbf{现时位置}与速度的夹角：
\[
\重点{\mathbf E=\frac{q}{4\pi\varepsilon_0}\frac{1-\beta^2}{(1-\beta^2\sin^2\theta)^{3/2}}\frac{\hat{\mathbf R}}{R^2}},\qquad \mathbf B=\frac{\mathbf v\times\mathbf E}{c^2}.
\]
横向场增强 $\gamma$ 倍，纵向削弱 $\gamma^{-2}$ 倍；掠过距离 $b$ 的冲量 $\Delta p_\perp=\frac{2kqQ}{bv}$（$k=\frac1{4\pi\varepsilon_0}$）与 $\gamma$ 无关。
\end{结论}

\subsection{热力学与统计物理}

\begin{结论}{配分函数的四条导数}
$\beta=1/kT$：
\[
F=-kT\ln Z,\quad U=-\partial_\beta\ln Z,\quad \重点{\langle(\Delta E)^2\rangle=\partial_\beta^2\ln Z=kT^2C_V},\quad S=\frac{U-F}{T}.
\]
二能级（间隙 $\Delta$）：$C=Nk\,x^2\frac{\ee^{x}}{(1+\ee^x)^2}$，$x=\Delta/kT$，在 $kT\approx0.42\Delta$ 处出现 Schottky 峰。独立子系统 $Z=Z_1^N$（可分辨）或 $Z_1^N/N!$（不可分辨、经典极限）。
\end{结论}

\begin{结论}{Jacobi 行列式推 Maxwell 关系}
$\dfrac{\partial(T,S)}{\partial(p,V)}=1$，从而 $\Bigl(\dfrac{\partial S}{\partial V}\Bigr)_T=\Bigl(\dfrac{\partial p}{\partial T}\Bigr)_V$ 等四式。常用结果：
\[
\重点{C_p-C_V=\frac{TV\alpha^2}{\kappa_T}},\qquad \Bigl(\frac{\partial U}{\partial V}\Bigr)_T=T\Bigl(\frac{\partial p}{\partial T}\Bigr)_V-p,\qquad \mu_{\mathrm{JT}}=\frac{V(T\alpha-1)}{C_p}.
\]
$\alpha=\frac1V(\partial_TV)_p$，$\kappa_T=-\frac1V(\partial_pV)_T$。理想气体 $T\alpha=1$，Joule--Thomson 系数为零。
\end{结论}

\begin{结论}{简并 Fermi 气体与 Bose 凝聚}
自旋 $\tfrac12$ 自由电子，数密度 $n$：
\[
\重点{E_F=\frac{\hbar^2}{2m}(3\pi^2n)^{2/3}},\qquad C_V=\frac{\pi^2}{2}Nk\frac{T}{T_F},\qquad P=\frac25nE_F\ (T=0).
\]
理想 Bose 气体（无自旋）：\重点{kT_c=\frac{2\pi\hbar^2}{m}\Bigl(\frac{n}{\zeta(3/2)}\Bigr)^{2/3}}，$T<T_c$ 时凝聚比例 $1-(T/T_c)^{3/2}$。二维理想 Bose 气体无凝聚（积分发散）。
\end{结论}

\begin{结论}{黑体与 Debye：$T^4$ 与 $T^3$}
\[
\重点{u=\frac{\pi^2k^4}{15\hbar^3c^3}T^4},\quad \text{出射通量 }\sigma T^4=\frac c4u,\quad \lambda_{\max}T=2.898\times10^{-3}\,\mathrm{m\,K},\quad p=\frac u3 .
\]
光子数密度 $\propto T^3$，每个光子平均能量 $\approx2.70\,kT$。固体低温：\重点{C_V=\frac{12\pi^4}{5}Nk\Bigl(\frac{T}{\Theta_D}\Bigr)^3}；高温回到 Dulong--Petit $3Nk$。
\end{结论}

\begin{结论}{Landau 平均场：指数与比热跃变}
$f=a(T-T_c)m^2+bm^4$（$a,b>0$）：
\[
m^2=\frac{a(T_c-T)}{2b}\ (T<T_c),\qquad \重点{\Delta C=\frac{a^2T_c}{2b}}\ (\text{单位体积}).
\]
平均场指数 $\beta=\frac12$，$\gamma=1$，$\delta=3$，$\alpha=0$（跃变）。Ising 平均场 $kT_c=zJ$（$z$ 为配位数）；一维无有限温相变，平均场在 $d<4$ 时给错误指数。
\end{结论}

\subsection{量子力学}

\begin{结论}{对易子三件套：$[x,f(p)]$、Ehrenfest、位力}
\[
\重点{[x,f(p)]=\ii\hbar f'(p)},\quad [p,g(x)]=-\ii\hbar g'(x),\quad \frac{\dd}{\dd t}\langle Q\rangle=\frac{\ii}{\hbar}\langle[H,Q]\rangle+\langle\partial_tQ\rangle .
\]
取 $Q=\mathbf r\cdot\mathbf p$ 于定态：\重点{2\langle T\rangle=\langle\mathbf r\cdot\nabla V\rangle}。幂律势 $V\propto r^n$ 给 $2\langle T\rangle=n\langle V\rangle$；氢原子 $\langle T\rangle=-E_n$。
\end{结论}

\begin{结论}{谐振子：升降算符直接算矩阵元}
\[
\重点{x=\sqrt{\frac{\hbar}{2m\omega}}\,(a+a^\dagger)},\quad p=\ii\sqrt{\frac{m\hbar\omega}{2}}\,(a^\dagger-a),\quad a|n\rangle=\sqrt n|n-1\rangle,\ a^\dagger|n\rangle=\sqrt{n+1}|n+1\rangle .
\]
$\langle n|x^2|n\rangle=\frac{\hbar}{m\omega}(n+\tfrac12)$，$\langle n|x^4|n\rangle=\frac{3\hbar^2}{4m^2\omega^2}(2n^2+2n+1)$。相干态 $a|\alpha\rangle=\alpha|\alpha\rangle$，光子数为 Poisson 分布，均值 $|\alpha|^2$，$\Delta x\Delta p=\hbar/2$。
\end{结论}

\begin{结论}{非简并微扰：二阶公式与符号}
\[
E_n^{(1)}=V_{nn},\qquad \重点{E_n^{(2)}=\sum_{m\ne n}\frac{|V_{mn}|^2}{E_n^{(0)}-E_m^{(0)}}}.
\]
基态二阶修正恒 $\le0$。核对例：谐振子加 $V=\lambda x$，只有 $m=1$ 项，$E_0^{(2)}=-\frac{\lambda^2}{2m\omega^2}$，与配方的精确结果一致。简并能级先在简并子空间对角化 $V$，再用好基矢做非简并公式。
\end{结论}

\begin{结论}{Fermi 黄金规则}
常数微扰 $V$ 打开后，跃迁到连续态的速率：
\[
\重点{\Gamma_{i\to f}=\frac{2\pi}{\hbar}|\langle f|V|i\rangle|^2\rho(E_f)}.
\]
简谐微扰 $V\ee^{-\ii\omega t}+\text{h.c.}$ 时 $\rho$ 在 $E_f=E_i\pm\hbar\omega$ 取值。适用时间窗：$\hbar/\Delta E_{\text{连续带}}\ll t\ll 1/\Gamma$。
\end{结论}

\begin{结论}{角动量耦合：$\mathbf L\cdot\mathbf S$ 与 Landé 因子}
在 $|j,m_j\rangle$ 基下：
\[
\重点{\langle\mathbf L\cdot\mathbf S\rangle=\frac{\hbar^2}{2}\bigl[j(j+1)-l(l+1)-s(s+1)\bigr]},\qquad g_J=1+\frac{j(j+1)+s(s+1)-l(l+1)}{2j(j+1)}.
\]
弱场 Zeeman：$\Delta E=g_J\mu_Bm_jB$。$j_1\otimes j_2=\bigoplus_{j=|j_1-j_2|}^{j_1+j_2}j$，维数核对 $(2j_1+1)(2j_2+1)$。两个自旋 $\tfrac12$：$\mathbf S_1\cdot\mathbf S_2=\frac{\hbar^2}{4}$（三重态）或 $-\frac{3\hbar^2}{4}$（单态）。
\end{结论}

\begin{结论}{变分法：试探函数的标度参数}
任意归一化 $\psi$：$\langle H\rangle\ge E_0$。对含标度参数 $b$ 的试探族，$\langle T\rangle\propto b^{-2}$、$\langle V\rangle\propto b^{n}$，极小化给 \重点{2\langle T\rangle=n\langle V\rangle}（试探族内的位力定理）。
氢原子用 Gauss 试探 $\ee^{-r^2/b^2}$：$E_{\min}=-\frac{8}{3\pi}\,\mathrm{Ry}\approx-0.85\,\mathrm{Ry}$。与基态正交的试探函数给第一激发态的上界。
\end{结论}

\begin{结论}{WKB：量子化条件与隧穿}
\[
\重点{\oint p\,\dd x=2\pi\hbar\,(n+\tfrac12)}\ (\text{两个光滑转折点});\quad \text{一侧硬墙 }n+\tfrac34;\quad \text{两侧硬墙 }n+1 .
\]
势垒穿透 \重点{T\approx\exp\Bigl(-\frac2\hbar\int_{x_1}^{x_2}\sqrt{2m(V-E)}\,\dd x\Bigr)}。Gamow：$\alpha$ 衰变寿命的对数与 $Z/\sqrt{E}$ 近似线性（Geiger--Nuttall）。
\end{结论}

\begin{结论}{氢原子的径向期望值}
$a_0$ 为 Bohr 半径，$E_n=-\frac{13.6\ \mathrm{eV}}{n^2}$（$Z=1$；类氢 $\times Z^2$，长度 $\div Z$）：
\[
\重点{\Bigl\langle\frac1r\Bigr\rangle=\frac{1}{n^2a_0}},\quad \langle r\rangle=\frac{a_0}{2}\bigl[3n^2-l(l+1)\bigr],\quad \Bigl\langle\frac{1}{r^2}\Bigr\rangle=\frac{1}{n^3(l+\frac12)a_0^2}.
\]
$\langle 1/r^3\rangle=\frac{1}{n^3l(l+\frac12)(l+1)a_0^3}$（$l\ge1$）用于精细结构；$\langle1/r\rangle$ 也可由 Feynman--Hellmann 对 $e^2$ 求导得到。
\end{结论}

\参考{Goldstein, \emph{Classical Mechanics}；Landau--Lifshitz 第 1 卷（MIT 8.09 / 8.223）；Griffiths, \emph{Introduction to Electrodynamics}；Jackson, \emph{Classical Electrodynamics}（8.07）；Kardar, \emph{Statistical Physics of Particles}（8.333）；Schroeder, \emph{Thermal Physics}（8.044）；Sakurai, \emph{Modern Quantum Mechanics}；Griffiths, \emph{Introduction to Quantum Mechanics}（8.04--8.06）。}

\续{07b}
\续{07c}
\续{07d}
\续{07e}
\续{07f}
\续{07g}
\section{电磁学}

\subsection{核心积木}

\begin{结论}{导体平面与导体球的镜像}
平面：电荷 $q$ 距接地平面 $d$，感应面密度 $\sigma=-\dfrac{qd}{2\pi(\rho^2+d^2)^{3/2}}$（总量 $-q$），引力 $\dfrac{q^2}{16\pi\varepsilon_0d^2}$。
接地球（半径 $R$，$q$ 距球心 $d>R$）：
\[
\重点{q'=-\frac{R}{d}q,\qquad b=\frac{R^2}{d}}\ (\text{镜像在球心与 }q\text{ 连线上}).
\]
孤立不带电球再在球心加 $+\frac Rdq$；带电 $Q$ 的孤立球在球心加 $Q+\frac Rdq$。
\end{结论}

\begin{结论}{偶极场与球内平均场}
\[
\重点{\mathbf E_{\mathrm{dip}}=\frac{1}{4\pi\varepsilon_0}\frac{3(\mathbf p\cdot\hat{\mathbf r})\hat{\mathbf r}-\mathbf p}{r^3}},\quad \mathbf F=\nabla(\mathbf p\cdot\mathbf E),\quad \boldsymbol\tau=\mathbf p\times\mathbf E,\quad U=-\mathbf p\cdot\mathbf E .
\]
半径 $R$ 球内的平均场：电荷都在球外时 $=$ 球心处的场；都在球内时 $=-\dfrac{\mathbf p}{4\pi\varepsilon_0R^3}$。磁偶极同形（$\mathbf p\to\mathbf m$，$\frac1{\varepsilon_0}\to\mu_0$），球内平均 $B=\frac{2\mu_0}{3}\frac{\mathbf m}{\frac43\pi R^3}$。
\end{结论}

\begin{结论}{均匀外场中的球：一个偶极就够}
外场 $E_0\hat z$，球外 $V=-E_0r\cos\theta+\frac{A\cos\theta}{r^2}$，球内 $V=-E_{\mathrm{in}}r\cos\theta$。
\[
\text{导体球：}p=4\pi\varepsilon_0R^3E_0,\ \sigma=3\varepsilon_0E_0\cos\theta;\qquad \text{介质球：}\重点{\mathbf E_{\mathrm{in}}=\frac{3}{\varepsilon_r+2}\mathbf E_0}.
\]
均匀极化球内 $\mathbf E=-\frac{\mathbf P}{3\varepsilon_0}$；均匀磁化球内 $\mathbf B=\frac23\mu_0\mathbf M$。
\end{结论}

\begin{结论}{能量法求力：固定 $Q$ 与固定 $V$ 的符号}
\[
\重点{F_x=-\Bigl(\frac{\partial W}{\partial x}\Bigr)_Q=+\Bigl(\frac{\partial W}{\partial x}\Bigr)_V},\qquad W=\tfrac12CV^2=\frac{Q^2}{2C}.
\]
固定 $V$ 时电池做功 $V\dd Q=2\,\dd W$，这就是符号反转的原因。介质板（宽 $w$、间距 $d$）被吸入平行板：$F=\dfrac{\varepsilon_0(\varepsilon_r-1)wV^2}{2d}$，与插入深度无关。
\end{结论}

\begin{结论}{Biot--Savart 的两条基本结果}
有限直导线，角从场点到导线的垂线量起：$B=\dfrac{\mu_0I}{4\pi s}(\sin\theta_2-\sin\theta_1)$。圆环轴线：
\[
\重点{B_z=\frac{\mu_0IR^2}{2(R^2+z^2)^{3/2}}}.
\]
Helmholtz 线圈（两环间距 $=R$）使中点 $\partial_z^2B=0$，中心 $B=\bigl(\tfrac45\bigr)^{3/2}\frac{\mu_0NI}{R}$；正 $n$ 边形线圈中心场由 $n$ 段直导线叠加。
\end{结论}

\begin{结论}{同轴线：$L$、$C$ 与特征阻抗}
内半径 $a$、外半径 $b$，单位长度：
\[
\frac{L}{\ell}=\frac{\mu}{2\pi}\ln\frac ba,\quad \frac C\ell=\frac{2\pi\varepsilon}{\ln(b/a)},\quad \重点{Z_0=\frac{1}{2\pi}\sqrt{\frac\mu\varepsilon}\ln\frac ba},\quad \frac{L}{\ell}\cdot\frac C\ell=\mu\varepsilon .
\]
$LC=\mu\varepsilon$ 对任何横截面的 TEM 传输线都成立，信号速度 $1/\sqrt{\mu\varepsilon}$。（$L$ 为外导体内电流的外电感，忽略导体内部。）
\end{结论}

\begin{结论}{动生电动势：终端速度与电磁阻尼}
导轨间距 $\ell$，电阻 $R$，杆质量 $m$ 在重力下沿竖直导轨滑落，垂直磁场 $B$：
\[
m\dot v=mg-\frac{B^2\ell^2}{R}v,\qquad \重点{v_\infty=\frac{mgR}{B^2\ell^2}},\qquad \tau=\frac{mR}{B^2\ell^2}.
\]
磁通法则 $\mathcal E=-\dot\Phi$ 在回路形状改变时仍成立，但“导体在磁场中不动、只是接触点移动”的情形须回到 $\oint(\mathbf E+\mathbf v\times\mathbf B)\cdot\dd\boldsymbol\ell$。
\end{结论}

\begin{结论}{RLC：特征量与充电损耗}
串联 RLC：$\omega_0=\frac{1}{\sqrt{LC}}$，\重点{Q=\frac1R\sqrt{\frac LC}}，临界阻尼 $R=2\sqrt{L/C}$；时间常数 $RC$、$L/R$。
恒压 $V$ 经任意电阻给电容充满：电池做功 $CV^2$，电容储能 $\frac12CV^2$，\重点{\text{电阻耗散 }\tfrac12CV^2\text{，与 }R\text{ 无关}}。两等容电容一充一空相连，同样损失一半能量。
\end{结论}

\begin{结论}{等效电阻的两个对称技巧}
无限方格电阻网（每边 $R$），相邻两节点间 \重点{R_{\mathrm{eq}}=R/2}：注入 $I$ 与抽出 $I$ 两种情形叠加，每种情形该边分到 $I/4$。
无限梯形网（串 $R$、并 $R$）：自相似 $R_\infty=R+\frac{RR_\infty}{R+R_\infty}$，得 $R_\infty=\frac{1+\sqrt5}{2}R$。立方体对顶点间 $\frac56R$。
\end{结论}

\begin{结论}{界面条件与场线折射}
\[
\重点{D_{1\perp}-D_{2\perp}=\sigma_f,\quad \mathbf E_{1\parallel}=\mathbf E_{2\parallel},\quad B_{1\perp}=B_{2\perp},\quad \hat{\mathbf n}\times(\mathbf H_1-\mathbf H_2)=\mathbf K_f}.
\]
无自由面电荷时 $\dfrac{\tan\theta_1}{\tan\theta_2}=\dfrac{\varepsilon_1}{\varepsilon_2}$（$\theta$ 从法线量）；高 $\mu$ 材料内磁感线几乎平行于表面，这是磁屏蔽的原理。
\end{结论}

\begin{结论}{趋肤深度与电荷弛豫}
良导体（$\sigma\gg\varepsilon\omega$）：\重点{\delta=\sqrt{\frac{2}{\mu\sigma\omega}}}；铜在 $60\,$Hz 约 $8.4\,$mm，$1\,$MHz 约 $65\,\mu$m。导体内自由电荷按 $\ee^{-t/\tau}$ 消散，\重点{\tau=\varepsilon/\sigma}。
良导体反射：表面阻抗 $Z_s=(1+\ii)/(\sigma\delta)$，正入射吸收比例 $A\approx\dfrac{4\,\Re Z_s}{Z_0}=\dfrac{2\omega\delta}{c}$（$\mu=\mu_0$，领头阶）。
\end{结论}

\begin{结论}{电路中的能量走 Poynting 矢量}
载流 $I$ 的直导线（电阻 $R$，长 $\ell$，半径 $a$）：表面 $E=IR/\ell$，$B=\frac{\mu_0I}{2\pi a}$，
\[
\重点{\oint\mathbf S\cdot\dd\mathbf a=\frac{EB}{\mu_0}\cdot2\pi a\ell=I^2R}.
\]
能量从外部场流入导线而非沿导线内部流动；充电电容中能量同样从侧面流入极板间隙。
\end{结论}

\参考{Purcell--Morin, \emph{Electricity and Magnetism}（MIT 8.022）；Griffiths, \emph{Introduction to Electrodynamics}（8.02 / 8.07）；Zangwill, \emph{Modern Electrodynamics}；Horowitz--Hill, \emph{The Art of Electronics}。}

\续{08b}
\续{08c}
\续{08d}
\续{08e}
\section{近代物理}

\subsection{核心积木}

\begin{结论}{阈能：在质心系算不变质量}
$s=(\sum E)^2-(\sum\mathbf pc)^2$ 是不变量。束流 $m_1$ 打静止靶 $m_2$，产生总静质量 $\sum m_f$：
\[
\重点{E_{1,\mathrm{th}}=\frac{(\sum m_f)^2-m_1^2-m_2^2}{2m_2}\,c^2}\quad(\text{束流总能量}).
\]
$pp\to ppp\bar p$：$E_{\mathrm{th}}=7m_pc^2$，动能 $6m_pc^2\approx5.6\,$GeV。对撞机对头碰撞 $\sqrt s=2E$，固定靶 $\sqrt s\approx\sqrt{2m_2c^2E}$，这就是对撞机的优势。
\end{结论}

\begin{结论}{二体衰变的能量与动量}
静止 $M\to m_1+m_2$：
\[
\重点{E_1=\frac{M^2+m_1^2-m_2^2}{2M}c^2},\qquad |\mathbf p|c=\frac{\sqrt{[M^2-(m_1+m_2)^2][M^2-(m_1-m_2)^2]}}{2M}c^2 .
\]
$\pi^+\to\mu^+\nu$：$p_\mu c\approx29.8\,$MeV，$T_\mu\approx4.1\,$MeV（单一谱线）；三体衰变如 $\beta$ 衰变则为连续谱，这正是中微子假说的来源。
\end{结论}

\begin{结论}{Compton 散射}
\[
\重点{\lambda'-\lambda=\frac{h}{m_ec}(1-\cos\theta)},\qquad \lambda_C=\frac{h}{m_ec}=2.426\ \mathrm{pm},\qquad E'=\frac{E}{1+\frac{E}{m_ec^2}(1-\cos\theta)} .
\]
反冲电子最大动能（$\theta=\pi$）$T_{\max}=\frac{2E^2}{m_ec^2+2E}$，即 Compton 边。逆 Compton：极端相对论电子把低能光子能量提升约 $\frac43\gamma^2$ 倍（各向同性平均）。
\end{结论}

\begin{结论}{相对论 Doppler 与光行差}
源以 $\beta$ 运动，$\theta$ 为观察者系中速度与源到观察者方向的夹角：
\[
\重点{f=\frac{f_0}{\gamma(1-\beta\cos\theta)}},\qquad \text{正对接近 }f=f_0\sqrt{\frac{1+\beta}{1-\beta}},\qquad \cos\theta=\frac{\cos\theta'+\beta}{1+\beta\cos\theta'} .
\]
$\theta=\pi/2$ 时为横向红移 $f_0/\gamma$。源系各向同性发射的一半光子集中在观察者系半角 $\sim1/\gamma$ 的前向锥内（相对论束流）。
\end{结论}

\begin{结论}{快度与匀加速运动}
\重点{\tanh\eta=\beta}，共线速度合成即快度相加；$E=mc^2\cosh\eta$，$p=mc\sinh\eta$。固有加速度 $g$ 恒定（从静止出发）：
\[
\重点{t=\frac cg\sinh\frac{g\tau}{c}},\qquad x=\frac{c^2}{g}\Bigl(\cosh\frac{g\tau}c-1\Bigr).
\]
距离 $c^2/g$ 之后的光信号永远追不上（Rindler 视界）；$g=9.8\,\mathrm{m/s^2}$ 时 $c^2/g\approx0.97$ 光年。
\end{结论}

\begin{结论}{类氢体系的约化质量标度}
\[
\重点{E_n=-\frac{\mu}{m_e}\frac{Z^2}{n^2}\,13.6\ \mathrm{eV}},\qquad r_n=\frac{m_e}{\mu}\frac{n^2}{Z}a_0 .
\]
电子偶素 $\mu=m_e/2$：基态 $-6.8\,$eV；$\mu$ 子氢 $\mu\approx186m_e$：半径缩小约 186 倍。Moseley：$K_\alpha$ 能量 $\approx\frac34(Z-1)^2\times13.6\,$eV（$Z-1$ 为另一个 $1s$ 电子的屏蔽）。
\end{结论}

\begin{结论}{氢原子精细结构与超精细}
\[
\重点{E_{nj}=-\frac{13.6\ \mathrm{eV}}{n^2}\Bigl[1+\frac{\alpha^2}{n^2}\Bigl(\frac{n}{j+\frac12}-\frac34\Bigr)\Bigr]},\qquad \alpha\approx\frac1{137}.
\]
只依赖 $n,j$：$2s_{1/2}$ 与 $2p_{1/2}$ 简并，被 Lamb 位移（QED，约 $1057\,$MHz）打破。基态超精细分裂 $5.9\,\mu$eV，即 $21\,$cm（$1420\,$MHz）谱线；$\Delta E_{\mathrm{hf}}=\frac83g_p\frac{m_e}{m_p}\alpha^2\,\mathrm{Ry}$，$g_p\approx5.59$。
\end{结论}

\begin{结论}{常用数值组合}
\[
\重点{hc=1240\ \mathrm{eV\,nm}},\quad \hbar c=197.3\ \mathrm{MeV\,fm},\quad e^2/(4\pi\varepsilon_0)=1.44\ \mathrm{MeV\,fm},\quad m_ec^2=0.511\ \mathrm{MeV}.
\]
非相对论电子经电压 $V$ 加速：\重点{\lambda=\frac{1.226\ \mathrm{nm}}{\sqrt{V/\mathrm{V}}}}；热中子（$kT\approx25\,$meV）$\lambda\approx0.18\,$nm，适合晶体衍射。$\hbar=6.58\times10^{-22}\,$MeV\,s 把宽度换寿命：$\tau=\hbar/\Gamma$。
\end{结论}

\begin{结论}{Rutherford 散射}
入射动能 $E_k$，电荷 $ze$ 打 $Ze$，最近距离（正碰）$d=\frac{zZe^2}{4\pi\varepsilon_0E_k}$：
\[
\重点{b=\frac d2\cot\frac\theta2},\qquad \frac{\dd\sigma}{\dd\Omega}=\Bigl(\frac{d}{4}\Bigr)^2\frac{1}{\sin^4(\theta/2)} .
\]
$\sigma(\theta>\theta_0)=\pi b(\theta_0)^2$。厚 $x$、数密度 $n$ 的靶：散射比例 $n\,x\,\sigma$（薄靶）；衰减 $I=I_0\ee^{-n\sigma x}$，平均自由程 $1/(n\sigma)$。
\end{结论}

\begin{结论}{半经验质量公式与 $\beta$ 稳定线}
\[
B=a_VA-a_SA^{2/3}-a_C\frac{Z(Z-1)}{A^{1/3}}-a_A\frac{(A-2Z)^2}{A}+\delta,
\]
典型 $a_V\approx15.8$，$a_S\approx18.3$，$a_C\approx0.714$，$a_A\approx23.2\,$MeV。固定 $A$ 对 $Z$ 求原子质量极小（计入 $n$--$p$ 质量差）得 \重点{Z_{\mathrm{stable}}\approx\frac{A}{1.98+0.0155A^{2/3}}}；只用 $B$ 的近似为 $\frac{A}{2+(a_C/2a_A)A^{2/3}}$。$B/A$ 在 $^{56}$Fe 附近最大（约 $8.8\,$MeV），轻核聚变、重核裂变放能；核半径 $R\approx1.2A^{1/3}\,$fm。
\end{结论}

\begin{结论}{衰变链（Bateman 两步）}
$1\to2\to3$，$N_2(0)=0$：
\[
\重点{N_2(t)=\frac{\lambda_1N_1(0)}{\lambda_2-\lambda_1}\bigl(\ee^{-\lambda_1t}-\ee^{-\lambda_2t}\bigr)},\qquad t_{\max}=\frac{\ln(\lambda_2/\lambda_1)}{\lambda_2-\lambda_1}.
\]
$\lambda_1\ll\lambda_2$：长期平衡，活度 $\lambda_1N_1=\lambda_2N_2$；$\lambda_1<\lambda_2$ 但同量级：暂时平衡，$\frac{A_2}{A_1}\to\frac{\lambda_2}{\lambda_2-\lambda_1}$。
\end{结论}

\begin{结论}{电偶极跃迁选择定则}
\[
\重点{\Delta l=\pm1,\quad \Delta m_l=0,\pm1,\quad \Delta S=0,\quad \text{宇称改变}};\qquad \Delta J=0,\pm1\ (J=0\not\to0).
\]
速率 $\Gamma=\frac{\omega^3|\langle f|e\mathbf r|i\rangle|^2}{3\pi\varepsilon_0\hbar c^3}$；氢 $2p\to1s$ 寿命 $1.6\,$ns，$2s$ 只能双光子衰变（约 $0.12\,$s）。$\Delta m=0$ 为 $\pi$ 光（沿场线偏振），$\pm1$ 为 $\sigma^\pm$。
\end{结论}

\begin{结论}{弱场广义相对论的四个数}
\[
\重点{\frac{\Delta f}{f}=\frac{gh}{c^2}},\quad \delta\phi_{\text{光偏折}}=\frac{4GM}{c^2b},\quad \Delta\varpi=\frac{6\pi GM}{c^2a(1-e^2)}\ (\text{每圈}),\quad r_s=\frac{2GM}{c^2}.
\]
太阳边缘偏折 $1.75''$，水星进动每世纪 $43''$。GPS 卫星钟：狭义相对论 $-7\,\mu$s/天，引力 $+45\,\mu$s/天，净快约 $38\,\mu$s/天。
\end{结论}

\参考{French, \emph{Special Relativity}；Taylor--Wheeler, \emph{Spacetime Physics}（MIT 8.033 / 8.20）；Krane, \emph{Introductory Nuclear Physics}；Griffiths, \emph{Introduction to Elementary Particles}；Foot, \emph{Atomic Physics}；PDG 常数表。}

\续{09b}
\续{09c}
\续{09d}
\续{09e}
\section{光学波动}

\subsection{核心积木}

\begin{结论}{Fresnel 公式、Brewster 角与倏逝波}
\[
r_s=\frac{n_1\cos\theta_i-n_2\cos\theta_t}{n_1\cos\theta_i+n_2\cos\theta_t},\qquad r_p=\frac{n_2\cos\theta_i-n_1\cos\theta_t}{n_2\cos\theta_i+n_1\cos\theta_t}.
\]
正入射 $R=\bigl(\frac{n_1-n_2}{n_1+n_2}\bigr)^2$（玻璃约 $4\%$）；\重点{\tan\theta_B=n_2/n_1} 时 $r_p=0$。全反射时透射场按 $\ee^{-z/\delta}$ 衰减，\重点{\delta=\frac{1}{k_0\sqrt{n_1^2\sin^2\theta_i-n_2^2}}}。（$r_p$ 的整体符号随约定而异，$|r_p|$ 不变。）
\end{结论}

\begin{结论}{薄膜干涉与四分之一波增透}
膜厚 $t$、折射率 $n_f$，两反射光程差 $2n_ft\cos\theta_t$，再按\textbf{从低到高折射率反射}的次数补 $\lambda/2$：
\[
\text{增透：}\ \重点{n_ft=\frac\lambda4,\quad n_f=\sqrt{n_1n_2}}\ \Rightarrow\ R=0\ (\text{设计波长、正入射}).
\]
楔形膜条纹间距 $\frac{\lambda}{2n\alpha}$；Newton 环暗环半径 $r_m=\sqrt{m\lambda R}$（反射光，中心暗）。
\end{结论}

\begin{结论}{Fabry--Pérot：精细度与分辨率}
\[
T=\frac{1}{1+F\sin^2(\delta/2)},\quad \delta=\frac{4\pi nd\cos\theta}{\lambda},\quad F=\frac{4R}{(1-R)^2},\quad \重点{\mathcal F=\frac{\pi\sqrt R}{1-R}}.
\]
自由光谱区 $\Delta\nu_{\mathrm{FSR}}=\frac{c}{2nd}$，线宽 $\Delta\nu_{\mathrm{FSR}}/\mathcal F$，分辨本领 \重点{\frac{\lambda}{\delta\lambda}=m\mathcal F}。$R=0.99$ 时 $\mathcal F\approx313$。
\end{结论}

\begin{结论}{光栅：强度分布与分辨本领}
$N$ 缝、缝距 $d$、缝宽 $a$；$\beta=\frac{\pi a\sin\theta}{\lambda}$，$\alpha=\frac{\pi d\sin\theta}{\lambda}$：
\[
I=I_0\Bigl(\frac{\sin\beta}{\beta}\Bigr)^2\Bigl(\frac{\sin N\alpha}{\sin\alpha}\Bigr)^2,\qquad \重点{\frac{\lambda}{\Delta\lambda}=mN}.
\]
主极大 $d\sin\theta=m\lambda$，相邻主极大间有 $N-1$ 个零点；$d/a$ 为整数 $k$ 时 $m=k,2k,\dots$ 级缺级。角色散 $\frac{\dd\theta}{\dd\lambda}=\frac{m}{d\cos\theta}$。
\end{结论}

\begin{结论}{衍射极限与 Fresnel 数}
单缝暗纹 $a\sin\theta=m\lambda$；圆孔 Airy 第一暗环 \重点{\sin\theta=1.22\lambda/D}（Rayleigh 判据）。
\重点{N_F=\frac{a^2}{L\lambda}}：$N_F\ll1$ 为 Fraunhofer（远场 $=$ 孔径函数的 Fourier 变换，透镜焦面上精确成立），$N_F\gtrsim1$ 为 Fresnel。Babinet：互补屏在 Fraunhofer 区除零级外强度相同。
\end{结论}

\begin{结论}{Fresnel 半波带}
平面波照射，轴上距离 $L$，第 $m$ 带外半径 \重点{r_m=\sqrt{m\lambda L}}。每带贡献振幅近似等大、相邻反号，总振幅 $\approx$ 第一带的一半。
只露出第一带：强度为无遮挡时的 \重点{4} 倍；挡住奇数带的波带片焦距 $f=r_1^2/\lambda$（另有 $f/3,f/5,\dots$ 次焦点）。圆盘阴影中心总是亮点（Arago/Poisson 斑）。
\end{结论}

\begin{结论}{时间与空间相干}
相干时间 $\tau_c\approx\frac{1}{\Delta\nu}$，\重点{\ell_c=c\tau_c\approx\frac{\lambda^2}{\Delta\lambda}}：Michelson 光程差超过 $\ell_c$ 条纹消失。
van Cittert--Zernike：角直径 $\theta_s$ 的非相干源，横向相干宽度 \重点{\sim\lambda/\theta_s}（均匀圆盘源第一零点在 $1.22\lambda/\theta_s$）。双缝间距超过它条纹可见度消失——Michelson 恒星干涉仪测星角径的原理。
\end{结论}

\begin{结论}{Gauss 光束}
\[
w(z)=w_0\sqrt{1+(z/z_R)^2},\quad \重点{z_R=\frac{\pi w_0^2}{\lambda}},\quad R(z)=z\Bigl[1+\Bigl(\frac{z_R}{z}\Bigr)^2\Bigr],\quad \theta_{\mathrm{div}}=\frac{\lambda}{\pi w_0}.
\]
Gouy 相位 $\arctan(z/z_R)$，过焦点共移 $\pi$。焦距 $f$ 的透镜聚焦束腰 $w$ 的准直光：$w_0'\approx\frac{\lambda f}{\pi w}$。复参数 $q=z+\ii z_R$ 按 ABCD 律变换（见下条）。
\end{结论}

\begin{结论}{ABCD 矩阵与谐振腔稳定性}
自由传播 $\begin{pmatrix}1&d\\0&1\end{pmatrix}$，薄透镜 $\begin{pmatrix}1&0\\-1/f&1\end{pmatrix}$，曲率半径 $R$ 的反射镜同 $f=R/2$。成像条件 $B=0$，放大率 $A$。
\[
q'=\frac{Aq+B}{Cq+D};\qquad \text{两镜腔稳定：}\ \重点{0\le g_1g_2\le1,\quad g_i=1-\frac{L}{R_i}}.
\]
共焦腔 $g_1=g_2=0$，平行平面腔 $g_1=g_2=1$ 处于稳定边界。
\end{结论}

\begin{结论}{偏振：Jones 矩阵与 Malus}
快轴与 $x$ 轴成 $\theta$ 的半波片使线偏振方向从 $\phi$ 变为 \重点{2\theta-\phi}；快轴 $45^\circ$ 的四分之一波片把 $x$ 线偏振变为圆偏振。
Malus：$I=I_0\cos^2\theta$。自然光经三片依次转 $45^\circ$ 的偏振片：$\frac{I_0}{2}\cdot\frac12\cdot\frac12=\frac{I_0}{8}$；若去掉中间一片则为零。首片之后再放 $N$ 片、每片转 $\frac{\pi}{2N}$：$\frac{I_0}2\cos^{2N}\frac{\pi}{2N}\to\frac{I_0}{2}$（量子 Zeno 效应的光学类比）。
\end{结论}

\begin{结论}{色散：群速度与脉冲展宽}
\[
\重点{v_g=\frac{\dd\omega}{\dd k}=\frac{c}{n+\omega\frac{\dd n}{\dd\omega}}=\frac{c}{n-\lambda\frac{\dd n}{\dd\lambda}}},\qquad \beta_2=\frac{\dd^2k}{\dd\omega^2}.
\]
Gauss 脉冲（$1/\ee$ 强度半宽 $\tau_0$）：$\tau(z)=\tau_0\sqrt{1+(z/L_D)^2}$，$L_D=\tau_0^2/|\beta_2|$。脉冲越短，展宽越快（$L_D\propto\tau_0^2$）。
\end{结论}

\begin{结论}{阻抗失配：反射、透射与匹配}
弦 $Z=\sqrt{T\mu}$，声 $Z=\rho v$，传输线 $Z_0$。振幅（位移/电压）：
\[
\重点{r=\frac{Z_1-Z_2}{Z_1+Z_2}},\quad t=\frac{2Z_1}{Z_1+Z_2},\quad R+T=1\ \text{其中 }T=\frac{4Z_1Z_2}{(Z_1+Z_2)^2}.
\]
电压反射系数 $\Gamma=\frac{Z_L-Z_0}{Z_L+Z_0}$（与弦的符号约定相反）。四分之一波段匹配：$Z_m=\sqrt{Z_1Z_2}$，与增透膜同一结构。
\end{结论}

\begin{结论}{波导截止与 $v_gv_p=c^2$}
矩形波导 $a\times b$（$a>b$）：
\[
\重点{\omega_{mn}=c\pi\sqrt{(m/a)^2+(n/b)^2}},\qquad k=\frac1c\sqrt{\omega^2-\omega_{mn}^2},\qquad v_gv_p=c^2 .
\]
主模 TE$_{10}$，截止波长 $2a$；TM 模要求 $m,n\ge1$。低于截止频率时场倏逝衰减。
\end{结论}

\begin{结论}{等离子体频率}
\[
\重点{\omega_p^2=\frac{ne^2}{\varepsilon_0m_e}},\qquad \omega^2=\omega_p^2+c^2k^2,\qquad f_p\approx9\sqrt{n/\mathrm{m^{-3}}}\ \mathrm{Hz}.
\]
$\omega<\omega_p$ 的电磁波被反射（电离层反射短波；金属对可见光反射、紫外透明）。$n(\omega)=\sqrt{1-\omega_p^2/\omega^2}<1$，但 $v_g<c$。
\end{结论}

\begin{结论}{Doppler、Mach 锥与 Cherenkov 角}
声波（介质静止，速度 $v$）：$f'=f\,\dfrac{v+v_o}{v-v_s}$（相互接近为正）。
超音速源：Mach 锥半角 $\sin\theta=v/v_s$；带电粒子在介质中 $v>c/n$：\重点{\cos\theta_C=\frac{1}{n\beta}}，阈值 $\beta>1/n$，单位长度光子数 $\propto\bigl(1-\frac{1}{n^2\beta^2}\bigr)$。
\end{结论}

\参考{Hecht, \emph{Optics}；Born--Wolf, \emph{Principles of Optics}；French, \emph{Vibrations and Waves}（MIT 8.03）；Saleh--Teich, \emph{Fundamentals of Photonics}；Siegman, \emph{Lasers}。}

\续{10b}
\续{10c}
\续{10d}
\续{10e}
\section{电磁实验}

\subsection{核心积木}

\begin{结论}{荷质比：Helmholtz 线圈 $+$ 电子束}
加速电压 $V$，圆轨道半径 $r$，线圈 $N$ 匝、半径 $R$、电流 $I$：
\[
\重点{\frac em=\frac{2V}{B^2r^2}},\qquad B=\Bigl(\frac45\Bigr)^{3/2}\frac{\mu_0NI}{R}.
\]
$r$ 以平方进入，是主要误差源：固定 $I$ 改变 $V$，拟合直线 $r^2=\frac{2m}{eB^2}V$，由斜率求 $e/m$，同时可检验截距（接触电势、地磁场）。地磁场分量应与线圈轴对齐后扣除，或正反电流平均。
\end{结论}

\begin{结论}{Millikan 油滴：两次终端速度}
无场下落 $v_1$，加场 $E$ 上升 $v_2$，油密度 $\rho$（扣除空气浮力），粘滞系数 $\eta$：
\[
a=\sqrt{\frac{9\eta v_1}{2\rho g}},\qquad \重点{q=\frac{6\pi\eta a\,(v_1+v_2)}{E}}.
\]
$a\sim1\,\mu$m 与空气平均自由程可比，用 Cunningham 修正 $\eta_{\mathrm{eff}}=\eta/(1+\frac{b}{pa})$。同一油滴改变电荷（电离源）后 $v_2$ 的跳变比 $q$ 本身更能直接显示量子化。
\end{结论}

\begin{结论}{Hall 效应与 van der Pauw 法}
厚 $t$ 的样品，电流 $I$，垂直磁场 $B$：
\[
\重点{V_H=\frac{IB}{nqt}},\qquad R_H=\frac1{nq}\ (\text{符号定载流子类型}),\qquad \mu=|R_H|\sigma .
\]
电极错位带来的不等势电压不随 $B$ 反号：取 $(\pm B,\pm I)$ 四种组合平均消去。任意形状薄片的方块电阻：$\ee^{-\pi R_A/R_s}+\ee^{-\pi R_B/R_s}=1$，对称时 \重点{R_s=\frac{\pi R}{\ln2}}。
\end{结论}

\begin{结论}{Johnson 噪声与散粒噪声：测 $k$ 与 $e$}
\[
\重点{\langle V^2\rangle=4kTR\,\Delta f},\qquad \重点{\langle\delta I^2\rangle=2eI\,\Delta f}.
\]
$300\,$K、$1\,$k$\Omega$：$4.07\,\mathrm{nV}/\sqrt{\mathrm{Hz}}$。测 $k$：改变 $R$ 或 $T$，做 $\langle V^2\rangle$ 对 $R$ 的直线，截距为放大器自身噪声；$\Delta f$ 必须用实测增益平方对频率的积分（等效噪声带宽），而非标称带宽。
\end{结论}

\begin{结论}{锁相放大器：乘法 $+$ 低通}
信号 $V_s\sin(\omega t+\phi)$ 乘参考 $V_r\sin\omega t$ 再低通：
\[
\重点{X=\tfrac12V_sV_r\cos\phi},\quad Y=\tfrac12V_sV_r\sin\phi,\quad R=\sqrt{X^2+Y^2}\ (\text{与相位无关}).
\]
时间常数 $\tau$ 的一阶 RC 低通等效噪声带宽 $\frac1{4\tau}$（6\,dB/oct），信噪比随 $\sqrt\tau$ 提高；调制频率选在 $1/f$ 噪声拐角之上且避开工频谐波。
\end{结论}

\begin{结论}{RC 滤波与运放增益}
一阶低通 $H=\frac{1}{1+\ii\omega RC}$：\重点{f_c=\frac{1}{2\pi RC}} 处 $-3\,$dB、相移 $-45^\circ$，之后 $-20\,$dB/十倍频。
理想运放（虚短、虚断）：反相 $-R_f/R_i$，同相 $1+R_f/R_i$，积分器 $-\frac{1}{RC}\int V\,\dd t$。增益带宽积 $GBW$ 固定：闭环增益 $G$ 时带宽约 $GBW/G$。
\end{结论}

\begin{结论}{电桥与四线法}
Wheatstone 平衡：\重点{R_x=R_3\frac{R_2}{R_1}}（与电源电压无关，零示法）。不平衡小量 $\delta$：输出 $\approx\frac{V}{4}\frac{\delta R}{R}$（四臂相等时）。
小电阻用四线法：电流线与电压线分开，电压表几乎无电流，引线与接触电阻不进入读数。交流电桥平衡须同时满足幅值与相位两个条件。
\end{结论}

\begin{结论}{Zeeman 效应：用 F--P 环读小分裂}
正常 Zeeman 分裂 $\Delta E=\mu_BB$，$\mu_B/h=14.0\,$GHz/T。F--P 标准具（间隔 $d$）的环满足 $r_p^2$ 等差，相邻环差 $\Delta(r^2)$；同一级内两分量的差 $\delta(r^2)$：
\[
\重点{\Delta\tilde\nu=\frac{1}{2d}\,\frac{\delta(r^2)}{\Delta(r^2)}}\quad(\tilde\nu=1/\lambda).
\]
只用平方半径的\textbf{比值}，焦距与级次都消掉。偏振：沿 $B$ 观察只见 $\sigma^\pm$（圆偏振），垂直观察见 $\pi$（平行 $B$）与 $\sigma$（垂直 $B$）。
\end{结论}

\begin{结论}{Franck--Hertz：峰间距而非峰位}
汞蒸气中电流峰间距 $\approx4.9\,$eV $=hc/253.7\,$nm（$6^3P_1$ 激发）。第一峰位置含接触电势，\重点{\text{只用相邻峰间距}}。
精细效应：电子达到激发阈值后还要再走约一个平均自由程 $\lambda$ 才碰撞，额外获得能量，使间距随级次近似线性增大（每级增量约 $E_a\lambda/L$，$L$ 为阴栅距离）；对级次外推到零可得更准的激发能。
\end{结论}

\begin{结论}{磁共振：Larmor 频率与自旋回波}
\重点{\omega=\gamma B}：质子 $\gamma/2\pi=42.58\,$MHz/T，自由电子 $g\mu_B/h=28.0\,$GHz/T。
$90^\circ$--$\tau$--$180^\circ$ 序列在 $2\tau$ 处出现回波，消去静态不均匀场：回波幅度 $\propto\ee^{-2\tau/T_2}$，而自由感应衰减按 $\frac1{T_2^*}=\frac1{T_2}+\gamma\Delta B$ 衰减。$T_1$ 用反转恢复 $M_z=M_0(1-2\ee^{-t/T_1})$，零点在 $t=T_1\ln2$。
\end{结论}

\begin{结论}{光电效应测 $h/e$}
\[
\重点{eV_s=hf-\phi}\ \Rightarrow\ V_s=\frac he f-\frac\phi e .
\]
对多条谱线做 $V_s$--$f$ 直线，斜率为 $h/e$，与光强无关。截止电压的读取受阳极反向电流与接触电势影响：后者只平移截距，不改斜率；前者需用“电流开始偏离本底”的判据统一处理。
\end{结论}

\begin{结论}{电子衍射：环半径到晶面间距}
加速电压 $V$，相对论修正：
\[
\重点{\lambda=\frac{h}{\sqrt{2m_eeV\bigl(1+\frac{eV}{2m_ec^2}\bigr)}}},\qquad 2d\sin\theta=\lambda,\qquad d\approx\frac{2L\lambda}{r}\ (\text{屏距 }L\text{，小角}).
\]
$10\,$kV 时修正约 $0.5\%$。多晶环按 $d_{hkl}=a/\sqrt{h^2+k^2+l^2}$ 指标化：比较 $r^2$ 之比与允许的 $h^2+k^2+l^2$ 序列（FCC：3,4,8,11,\dots）。
\end{结论}

\begin{结论}{谐振法测 $L$：把杂散电容线性化}
$f=\frac{1}{2\pi\sqrt{L(C+C_s)}}$ 中 $C_s$ 未知。改变已知 $C$，作
\[
\重点{\frac{1}{f^2}=4\pi^2L\,C+4\pi^2LC_s},
\]
斜率给 $L$，截距与斜率之比给 $C_s$。同一思路：未知偏移量只要以加性方式进入，就改写为线性关系由截距吸收。
\end{结论}

\参考{Melissinos--Napolitano, \emph{Experiments in Modern Physics}；MIT 8.13--8.14 Junior Lab 实验讲义；Horowitz--Hill, \emph{The Art of Electronics}；Bevington--Robinson, \emph{Data Reduction and Error Analysis}。}

\续{11b}
\续{11c}
\section{数据分析}

\subsection{核心积木}

\begin{结论}{误差传递：线性化与幂律速算}
\[
\sigma_f^2=\sum_i\Bigl(\frac{\partial f}{\partial x_i}\Bigr)^2\sigma_i^2+2\sum_{i<j}\frac{\partial f}{\partial x_i}\frac{\partial f}{\partial x_j}\Cov(x_i,x_j).
\]
\重点{f=\prod x_i^{a_i}\Rightarrow\Bigl(\frac{\sigma_f}{f}\Bigr)^2=\sum a_i^2\Bigl(\frac{\sigma_i}{x_i}\Bigr)^2}（独立）。只需改进 $|a_i|\sigma_i/x_i$ 最大的那一项；$1/3$ 以下的项平方后不足十分之一。$f$ 在误差范围内明显非线性时（如 $\sigma_x/x\gtrsim0.3$ 的 $1/x$），改用 Monte Carlo 抽样传递。
\end{结论}

\begin{结论}{加权平均与一致性检验}
\[
\重点{\bar x=\frac{\sum x_i/\sigma_i^2}{\sum1/\sigma_i^2},\qquad \sigma_{\bar x}^2=\frac{1}{\sum1/\sigma_i^2}},\qquad \chi^2=\sum\frac{(x_i-\bar x)^2}{\sigma_i^2}\sim\chi^2_{N-1}.
\]
$\chi^2/(N-1)$ 明显大于 1：误差被低估或有未知系统误差；PDG 做法是把 $\sigma_{\bar x}$ 乘 $S=\sqrt{\chi^2/(N-1)}$。\注意 共同的系统误差不会随 $N$ 减小，须在平均后单独加上。
\end{结论}

\begin{结论}{带误差的直线拟合：闭式解}
$y=a+bx$，权 $w_i=1/\sigma_i^2$；$S=\sum w$，$S_x=\sum wx$，$S_{xx}=\sum wx^2$，$S_y,S_{xy}$ 同理，$\Delta=SS_{xx}-S_x^2$：
\[
b=\frac{SS_{xy}-S_xS_y}{\Delta},\quad a=\frac{S_{xx}S_y-S_xS_{xy}}{\Delta},\quad \重点{\sigma_b^2=\frac{S}{\Delta},\ \sigma_a^2=\frac{S_{xx}}{\Delta},\ \Cov(a,b)=-\frac{S_x}{\Delta}}.
\]
以加权平均 $\bar x=S_x/S$ 为原点重写 $y=a'+b(x-\bar x)$，则 $a',b$ 不相关；外推到远离数据的 $x$ 时必须计入协方差项。
\end{结论}

\begin{结论}{两个变量都有误差：有效方差}
$y=f(x)$，$x$、$y$ 误差独立：把权换成
\[
\重点{\sigma_{\mathrm{eff},i}^2=\sigma_{y,i}^2+\bigl(f'(x_i)\bigr)^2\sigma_{x,i}^2},
\]
直线时为 $\sigma_y^2+b^2\sigma_x^2$，需对 $b$ 迭代一两次。不能简单交换 $x,y$ 的角色做普通回归，两者斜率估计不互为倒数。
\end{结论}

\begin{结论}{$\chi^2$ 的读法与 $\Delta\chi^2$ 置信域}
拟合优度：$\chi^2_{\min}/\nu\approx1\pm\sqrt{2/\nu}$（$\nu=N-p$）。过大：模型或误差有问题；过小：误差被高估或数据相关。
参数误差：\重点{\Delta\chi^2=1}（单参数，$68.3\%$，对其余参数取极小）；联合 2 参数 $68.3\%$ 区域为 $\Delta\chi^2=2.30$，$95.4\%$ 为 $6.18$（$95.0\%$ 为 $5.99$）。在极小点附近 \重点{C=2\,H^{-1}}，$H_{jk}=\partial_j\partial_k\chi^2$。
\end{结论}

\begin{结论}{计数实验：扣本底与最优时间分配}
Poisson：$\sigma_N=\sqrt N$。信号$+$本底计 $t_s$ 得 $N_s$，纯本底计 $t_b$ 得 $N_b$：
\[
r=\frac{N_s}{t_s}-\frac{N_b}{t_b},\qquad \sigma_r^2=\frac{N_s}{t_s^2}+\frac{N_b}{t_b^2},\qquad \重点{\frac{t_s}{t_b}=\sqrt{\frac{R_{s+b}}{R_b}}}\ (\text{总时间固定时最优}).
\]
计数少（$\lesssim10$）时不要用 $\sigma=\sqrt{N_i}$ 加权的 $\chi^2$（偏向低计数）；用 Poisson 似然 $-2\ln L=2\sum(\mu_i-N_i\ln\mu_i)+\text{常数}$。
\end{结论}

\begin{结论}{死时间修正}
真实率 $n$，观测率 $m$，死时间 $\tau$：
\[
\text{不可扩展型：}\ \重点{n=\frac{m}{1-m\tau}};\qquad \text{可扩展型：}\ m=n\,\ee^{-n\tau}\ (m\le\tfrac{1}{\ee\tau}).
\]
双源法测 $\tau$：$m_{12}<m_1+m_2-m_b$，不可扩展型、本底很小时一阶近似 $\tau\approx\frac{m_1+m_2-m_{12}-m_b}{2m_1m_2}$。可扩展型高计数时观测率反而下降，一个 $m$ 对应两个 $n$。
\end{结论}

\begin{结论}{线性化变换要连同权重一起变}
$y=Ax^k$ 取对数，$y=A\ee^{-E/kT}$ 作 $\ln y$ 对 $1/T$：变换后
\[
\重点{\sigma_{\ln y}=\frac{\sigma_y}{y}},\qquad \text{权}\ w_i=\frac{y_i^2}{\sigma_{y,i}^2}.
\]
等权拟合 $\ln y$ 会让小 $y$ 的点权重过大。截距 $\ln A$ 的误差换回 $A$ 时分布不对称：报 $A\,\ee^{\pm\sigma_{\ln A}}$。
\end{结论}

\begin{结论}{指数寿命的极大似然}
观测窗 $[0,\infty)$ 内 $n$ 个衰变时间：\重点{\hat\tau=\bar t}，$\sigma_{\hat\tau}=\hat\tau/\sqrt n$。只能测到 $[0,T]$ 时似然改为 $\prod\frac{\ee^{-t_i/\tau}}{\tau(1-\ee^{-T/\tau})}$，$\bar t$ 会低估 $\tau$。
直方图拟合（等宽 bin、各 bin 计数为 Poisson）用 Poisson 似然；若对 $\ln N$ 作直线拟合，空 bin 与小计数 bin 会产生偏差。
\end{结论}

\begin{结论}{小样本：$t$ 分布与“误差的误差”}
$n$ 次测量：$\bar x\pm t_{n-1}\frac{s}{\sqrt n}$，$95\%$ 时 $t_4=2.78$，$t_9=2.26$，$t_\infty=1.96$。
\[
\重点{\frac{\sigma_s}{\sigma}\approx\frac{1}{\sqrt{2(n-1)}}}:\quad n=10\ \text{时误差本身约有 }24\%\text{ 的不确定度}.
\]
故误差一般只报 1--2 位有效数字，中心值保留到与误差相同的数位。
\end{结论}

\begin{结论}{Jackknife 与 Bootstrap}
复杂统计量 $\hat\theta$（比值、拟合参数等）没有解析误差时：
\[
\重点{\widehat{\Var}_{\mathrm{jack}}=\frac{n-1}{n}\sum_{i=1}^n\bigl(\hat\theta_{(i)}-\bar\theta_{(\cdot)}\bigr)^2},\qquad \text{偏差修正 }\hat\theta-(n-1)(\bar\theta_{(\cdot)}-\hat\theta).
\]
$\hat\theta_{(i)}$ 为删去第 $i$ 个点后的估计。Bootstrap：有放回重抽 $n$ 个点重复 $B\sim10^3$ 次，用 $\hat\theta^*$ 的标准差或分位数。数据有时间相关时须按块重抽。
\end{结论}

\begin{结论}{采样与频谱：分辨率、Nyquist 与混叠}
采样率 $f_s$，总时长 $T=N/f_s$：频率分辨率 \重点{\Delta f=1/T}，最高可辨频率 $f_s/2$。频率 $f$ 的信号被记录为
\[
\重点{f_{\mathrm{alias}}=|f-kf_s|}\quad(\text{取 }k\text{ 使结果落在 }[0,f_s/2]).
\]
补零只是插值，不提高分辨率；非整周期截断产生泄漏，加 Hann 窗以主瓣加宽换旁瓣压低。采样前必须模拟低通滤去 $f_s/2$ 以上分量。
\end{结论}

\参考{Bevington--Robinson, \emph{Data Reduction and Error Analysis for the Physical Sciences}；Taylor, \emph{An Introduction to Error Analysis}；Hughes--Hase, \emph{Measurements and their Uncertainties}；Press et al., \emph{Numerical Recipes} 第 15 章；PDG 统计综述。}

\续{12b}
\续{12c}
\section{综合应用}

\subsection{核心积木}

\begin{结论}{Einstein 关系：扩散 $=$ 涨落 $\times$ 耗散}
迁移率 $\mu_m=v/F$：\重点{D=\mu_mkT}；球形粒子 Stokes 阻力给 $D=\dfrac{kT}{6\pi\eta a}$（Stokes--Einstein）。
每维 $\langle x^2\rangle=2Dt$。由 Brown 运动的 $\langle x^2\rangle$、$\eta$、$a$、$T$ 测出 $k$，再由 $R=N_Ak$ 得 Avogadro 数（Perrin）。离子：$D=\frac{\mu_ekT}{q}$（$\mu_e$ 为电迁移率）。
\end{结论}

\begin{结论}{Green--Kubo：扩散系数是速度相关的积分}
Langevin $m\dot v=-m\gamma v+\xi(t)$，$\langle\xi(t)\xi(t')\rangle=2m\gamma kT\,\delta(t-t')$：
\[
\langle v(0)v(t)\rangle=\frac{kT}{m}\ee^{-\gamma|t|},\qquad \重点{D=\int_0^\infty\langle v(0)v(t)\rangle\,\dd t=\frac{kT}{m\gamma}}.
\]
噪声强度由平衡时 $\frac12m\langle v^2\rangle=\frac12kT$ 唯一确定，这是涨落--耗散定理最简单的实例。
\end{结论}

\begin{结论}{光镊标定：均分定理与 Lorentz 谱}
过阻尼 $\gamma\dot x=-\kappa x+\xi$。两种独立标定刚度 $\kappa$：
\[
\重点{\kappa=\frac{kT}{\langle x^2\rangle}},\qquad S_x(f)=\frac{kT}{\pi^2\gamma\,(f^2+f_c^2)}\ (\text{单边}),\quad \重点{f_c=\frac{\kappa}{2\pi\gamma}}.
\]
谱对 $f$ 积分恰好回到 $kT/\kappa$。热噪声极限同样适用于 AFM 悬臂、扭秤：可测力的下限 $\sim\sqrt{4kT\gamma\,\Delta f}$。
\end{结论}

\begin{结论}{三个 Green 函数与 Born 近似}
\[
\nabla^2\frac{1}{4\pi r}=-\delta^3,\qquad (\nabla^2+k^2)\frac{\ee^{\ii kr}}{4\pi r}=-\delta^3,\qquad (\nabla^2-\mu^2)\frac{\ee^{-\mu r}}{4\pi r}=-\delta^3 .
\]
一阶 Born：\重点{f(\theta)=-\frac{m}{2\pi\hbar^2}\int\ee^{-\ii\mathbf q\cdot\mathbf r}V(\mathbf r)\,\dd^3r}，$q=2k\sin\frac\theta2$。Yukawa 势 $V=\beta\frac{\ee^{-\mu r}}{r}$ 给 $f=-\frac{2m\beta}{\hbar^2(q^2+\mu^2)}$；$\mu\to0$ 恰好重现 Rutherford 截面。Debye 屏蔽长度 $\lambda_D=\sqrt{\varepsilon_0kT/(ne^2)}$ 就是这里的 $1/\mu$。
\end{结论}

\begin{结论}{Kramers--Kronig：因果性约束实部与虚部}
$\chi(t)=0$（$t<0$）且 $\chi(\omega)$ 在上半平面解析、$\to0$：
\[
\重点{\chi'(\omega)=\frac2\pi\,\mathrm P\!\int_0^\infty\frac{\omega'\chi''(\omega')}{\omega'^2-\omega^2}\dd\omega'}.
\]
推论：$\chi'(0)=\frac2\pi\int_0^\infty\frac{\chi''(\omega')}{\omega'}\dd\omega'$——静态极化率由全部吸收决定；透明区的折射率色散来自别处的吸收线。
\end{结论}

\begin{结论}{两能级：避免交叉、Rabi 与 Landau--Zener}
$H=\begin{pmatrix}\Delta/2&V\\V&-\Delta/2\end{pmatrix}$，本征值 $\pm\sqrt{\Delta^2/4+V^2}$，交叉点能隙 $2|V|$。从态 1 出发：
\[
P_{1\to2}(t)=\frac{V^2}{V^2+\Delta^2/4}\sin^2\Bigl(\frac{\sqrt{V^2+\Delta^2/4}}{\hbar}t\Bigr),\qquad \重点{P_{\text{非绝热}}=\exp\Bigl(-\frac{2\pi V^2}{\hbar\,|\dot\Delta|}\Bigr)}.
\]
同一矩阵描述耦合摆的拍、NMR 的旋转坐标系、中微子振荡（$\Delta\propto\Delta m^2/E$）。
\end{结论}

\begin{结论}{Drude 模型与 Wiedemann--Franz}
\[
\sigma_0=\frac{ne^2\tau}{m},\qquad \sigma(\omega)=\frac{\sigma_0}{1-\ii\omega\tau},\qquad \重点{\frac{\kappa}{\sigma T}=\frac{\pi^2}{3}\Bigl(\frac ke\Bigr)^2=2.44\times10^{-8}\ \mathrm{W}\,\Omega\,\mathrm{K^{-2}}}.
\]
Lorenz 数用 Sommerfeld（Fermi--Dirac）比热得到；经典 Drude 的数值恰好接近，是两处因子 100 的误差互相抵消。$\omega\tau\gg1$ 时回到等离子体响应 $\varepsilon=1-\omega_p^2/\omega^2$。
\end{结论}

\begin{结论}{紧束缚能带：有效质量与 Bloch 振荡}
一维最近邻跃迁 $t$，晶格常数 $a$：
\[
\重点{E(k)=\varepsilon_0-2t\cos ka},\qquad v_g=\frac{2ta}{\hbar}\sin ka,\qquad m^*=\frac{\hbar^2}{2ta^2}\ (\text{带底}).
\]
带宽 $4t$；带顶 $m^*<0$（空穴）。恒定电场下 $\hbar\dot k=-eE$，无散射时做 Bloch 振荡，周期 $T_B=\frac{h}{eEa}$。
\end{结论}

\begin{结论}{Born--Oppenheimer 能标比}
电子 : 振动 : 转动 $\approx$ \重点{1:\sqrt{m_e/M}:m_e/M}。电子能标 $\sim\frac{\hbar^2}{m_ea_0^2}\sim$ 几 eV，振动 $\sim0.1\,$eV，转动 $\sim10^{-3}$--$10^{-4}\,$eV。
转动谱 $E_J=\frac{\hbar^2}{2I}J(J+1)$，相邻吸收线间距 $2B$（$B=\hbar^2/2I$）；室温双原子气体振动冻结、转动激发，$C_V\approx\frac52R$。
\end{结论}

\begin{结论}{最大熵 $=$ Lagrange 乘子}
在 $\sum p_i=1$、$\sum p_iE_i=U$ 约束下极大化 $-\sum p_i\ln p_i$：
\[
\重点{p_i=\frac{\ee^{-\beta E_i}}{Z}},\qquad \beta\ \text{即约束 }U\text{ 的乘子}.
\]
连续变量只约束均值与方差得 Gauss 分布；只约束正值均值得指数分布。同一推导给出统计力学正则系综和统计学中的指数族。
\end{结论}

\begin{结论}{Metropolis 与 Monte Carlo 误差}
接受率 \重点{\min\bigl(1,\ee^{-\beta\Delta E}\bigr)}（对称提议）满足细致平衡，故平稳分布为 Boltzmann 分布。
Monte Carlo 积分误差 $\sigma/\sqrt N$，与维数无关；Markov 链样本相关，有效样本数 $N/(2\tau_{\mathrm{int}})$，$\tau_{\mathrm{int}}=\frac12+\sum_{t\ge1}\rho(t)$。重要抽样把方差换成 $\Var_q(f p/q)$，$q\propto|fp|$ 时最小。
\end{结论}

\begin{结论}{辛积分器：为何长时间模拟用 Verlet}
谐振子 $\omega$，步长 $h$。前向 Euler 每步把能量乘 \重点{1+\omega^2h^2}，单调发散；
蛙跳/Verlet $x_{n+1}=2x_n-x_{n-1}-h^2\omega^2x_n$ 保辛，稳定条件 \重点{\omega h<2}，能量误差有界振荡而不漂移（守恒一个修正 Hamilton 量）。分子动力学与天体力学长时间积分的标准选择。
\end{结论}

\begin{结论}{时间--带宽积}
Gauss 脉冲（半高全宽）\重点{\Delta t\,\Delta\nu=0.441}；sech$^2$ 脉冲 $0.315$。乘积大于此值说明有啁啾。
同一 Fourier 关系：$\Delta x\,\Delta k\ge\frac12$（标准差），谱线自然宽度 $\Delta\nu=\frac{1}{2\pi\tau}$，测量时间 $T$ 的频率分辨率 $\sim1/T$。
\end{结论}

\begin{结论}{理想高分子链：熵弹簧}
$N$ 段、每段 $b$ 的自由连接链：$\langle R^2\rangle=Nb^2$，端距分布近似 Gauss，自由能 $F=\frac{3kTR^2}{2Nb^2}+\text{常数}$：
\[
\重点{f=\frac{3kT}{Nb^2}R}.
\]
弹性系数 $\propto T$：橡皮筋受热收缩（熵弹性）。自回避链 $R\sim N^{\nu}$，$\nu\approx0.588$（三维），Flory 估计 $\nu=\frac{3}{d+2}$。
\end{结论}

\参考{Reif, \emph{Fundamentals of Statistical and Thermal Physics}；Ashcroft--Mermin, \emph{Solid State Physics}；Sakurai, \emph{Modern Quantum Mechanics}；Newman--Barkema, \emph{Monte Carlo Methods in Statistical Physics}；Hairer--Lubich--Wanner, \emph{Geometric Numerical Integration}；Phillips et al., \emph{Physical Biology of the Cell}。}

\续{13b}
\续{13c}
\end{document}
